% concatenated from arxiv.tex
\documentclass{article}

\usepackage[final]{neurips_2025}

\usepackage[utf8]{inputenc}
\usepackage[T1]{fontenc}
\usepackage{hyperref}
\usepackage{url}
\usepackage{booktabs}
\usepackage{amsfonts}
\usepackage{nicefrac}
\usepackage{microtype}
\usepackage{xcolor}

\usepackage{amsmath}
\usepackage{amsthm}
\usepackage{amssymb}
\usepackage[pdftex]{graphicx}
\usepackage{complexity}
\usepackage{algorithm}
\usepackage[noend]{algpseudocode}
\algrenewcommand\algorithmicrequire{\textbf{Input:}}
\algrenewcommand\algorithmicensure{\textbf{Output:}}
\usepackage{cleveref}

\newtheorem{theorem}{Theorem}[section]
\newtheorem{definition}{Definition}
\newtheorem{corollary}[theorem]{Corollary}
\newtheorem{prop}[theorem]{Proposition}
\newtheorem{lemma}[theorem]{Lemma}
\newtheorem{claim}[theorem]{Claim}
\newtheorem{example}[theorem]{Example}
\newtheorem{fact}[theorem]{Fact}

\DeclareMathOperator{\sign}{sign}
\DeclareMathOperator{\conv}{conv}


\title{The Computational Complexity of Counting Linear Regions in ReLU Neural Networks}

\author{
  Moritz Stargalla \\
  University of Technology Nuremberg \\
  \texttt{moritz.stargalla@utn.de} \\
  \And
  Christoph Hertrich \\
  University of Technology Nuremberg \\
  \texttt{christoph.hertrich@utn.de} \\
  \And
  Daniel Reichman \\
  Worcester Polytechnic Institute \\
  \texttt{daniel.reichman@gmail.com}
}

\renewcommand{\paragraph}[1]{\textbf{#1}.}

\begin{document}


\maketitle


\begin{abstract}
  An established measure of the expressive power of a given ReLU neural network is the number of linear regions into which it partitions the input space.
  There exist many different, non-equivalent definitions of what a linear region actually is. We systematically assess which papers use which definitions and discuss how they relate to each other. We then analyze the computational complexity of counting the number of such regions for the various definitions. Generally, this turns out to be an intractable problem. We prove \NP{}- and \#\P{}-hardness results already for networks with one hidden layer and strong hardness of approximation results for two or more hidden layers. Finally, on the algorithmic side, we demonstrate that counting linear regions can at least be achieved in polynomial \emph{space} for some common definitions.
\end{abstract}


\section{Introduction}

Neural networks with rectified linear unit (ReLU) activations are among the most common and fundamental models in modern machine learning. The functions represented by ReLU networks are \emph{continuous and piecewise linear} (CPWL), meaning that the input space can be partitioned into finitely many pieces on each of which the function is affine. Such pieces are called \emph{linear regions}. This leads to the following intuition: the more linear regions a neural network can produce, the more complex problems it is capable of solving. Consequently, starting with \citet{pascanu2013number} and \citet{montufar2014number}, the number of linear regions became a standard measure of the expressive power of a ReLU network. Substantial effort has been put into understanding this quantity, e.g., by deriving upper and lower bounds depending on the network architecture or by developing algorithms to count it. More information can be found in the surveys \citet{huchette2023deep,balestriero2025geometry}.

Despite the significant interest in understanding the number of linear regions, surprisingly little is known about the most natural associated computational complexity question: Given a neural network,
what are the time and space requirements needed to determine how many regions it has? The main objective of our paper is to make progress on this question by proving 
complexity-theoretical results on the problem of counting linear regions.

However, before one can even talk about counting linear regions, one has to properly define them. What sounds like a simple exercise is actually a non-trivial task. In the literature, there exists a variety of non-equivalent definitions of what counts as a linear region of a ReLU network. For example, some authors define it via possible sets of active neurons, others define it solely based on the function represented by the neural network. Some authors require regions to be full-dimensional, or connected, or even convex, others do not.
Inconsistencies between definitions have led to confusion and even minor flaws in previous work, as we explain in \Cref{sec:flaws}.

\subsection{Our Contributions}\label{sec:contribution}

\paragraph{Definitions of linear regions}
In order to raise awareness to the technical, but important and non-trivial inconsistencies regarding the definition of linear regions in neural networks, we identify six non-equivalent, commonly used definitions in \Cref{sec:definitions}. We discuss how they relate to each other and provide a table demonstrating which authors used which definitions in previous work. We do not make a recommendation about what definition is the most reasonable one to use, as this depends on the context, but we encourage all authors of future papers to be aware of the subtleties carried by the different options and to be explicit about which definition they use and why. 

\paragraph{Complexity of counting regions in shallow networks} As for many questions regarding ReLU networks, it makes sense to first understand the most basic case with one hidden layer. In \Cref{sec:shallow}, we prove that, regardless of which of the six definitions one uses, the seemingly simple question of deciding whether a shallow network has more than one linear region can indeed be decided in polynomial time. 
However, for all six definitions, we show that determining the exact number of regions is \#\P-hard, meaning that, unless the commonly believed conjecture $\#\P\,\neq\,\FP$ fails, one cannot count regions of a shallow network in polynomial time. 
Furthermore, our reduction shows that even finding an algorithm that approximately counts the number of regions for one hidden layer might be intractable, as it would resolve long-standing open questions in the context of counting cells of hyperplane arrangements \citep{linial1986hard}.

\paragraph{Complexity of counting regions in networks with more than one hidden layer}
\citet{wang2022estimation} showed that deciding if a deep neural network has more than $K$ regions is NP-hard.
In \Cref{sec:deep} we improve upon \citet{wang2022estimation} in several aspects. While the hardness by \citet{wang2022estimation} only applies to networks with logarithmically growing depth (in the input dimension), we show that hardness can be proved for every constant number of hidden layers $\geq 2$ and even in the case $K=1$, that is, for deciding if the network has more than one linear region. Our reduction also implies running-time lower bounds based on the exponential-time hypothesis. We furthermore show that, unless common complexity assumptions fail, one cannot even approximate the number of regions within an exponential factor in polynomial time.

\paragraph{Counting regions using polynomial space}
While most of our results are concerned with lower bounds, in \Cref{sec:pspace}, we turn our attention towards proving an \emph{upper} bound on the computational complexity of region counting. \citet{wang2022estimation} proved\footnote{The proof by \citet{wang2022estimation} works for a different definition than claimed in their paper; see \Cref{sec:flaws}.} that for one definition of linear regions, the problem can be solved in exponential time. We show the stronger statement that for three of our definitions, polynomial space is sufficient.

\paragraph{Limitations}
Our paper is of theoretical nature and we strive towards a thorough understanding of the problem of counting regions from a computational complexity perspective.
As such, we naturally do not optimize our algorithms and reductions for efficiency or practical use, in contrast to, e.g., \citet{serra2018bounding} and \citet{cai2023getting}. Our hardness results are of worst-case nature. Consequently, although beyond the scope of our paper, it is conceivable that additional assumptions render the problem tractable.
For example, it would be very interesting to devise algorithms for region counting on networks that have been trained using gradient descent, as there is evidence that such networks have fewer regions~\citep{hanin2019deep}, which might allow faster algorithms.
Not all of our results are valid for all of the six definitions we identify. We discuss the open problems resulting from this in the context of the respective sections. In our list of definitions in \Cref{sec:definitions} and the corresponding \Cref{tab:literature}, we tried to capture the most relevant previous works on linear regions, but a full literature review, like \citet{huchette2023deep}, is beyond the scope of our paper.


\subsection{Related work}

\citet{huchette2023deep} survey polyhedral methods for deep learning, also treating the study of linear regions in detail. To the best of our knowledge, the first bounds on the number of regions in terms of the network architecture (e.g., number of neurons, network depth) were developed by \citet{pascanu2013number} and \citet{montufar2014number}. Subsequently, better bounds were established \citep{raghu2017expressive,arora2018understanding,serra2018bounding,zanotti2025bounds}. \citet{arora2018understanding} prove that every CPWL function can be represented by a ReLU network.

Several works have developed algorithms for enumerating linear regions.
\citet{serra2018bounding} and \citet{cai2023getting} present mixed-integer programming based routines to count the number of regions and \citet{masden2022algorithmic} presents an algorithm to enumerate the full combinatorial structure of activation regions. As discussed above, \citet{wang2022estimation} provides some initial results on the computational complexity of counting regions, which we strengthen significantly in this paper. Our reductions are related to other decision problems on trained neural networks, e.g., verification \citep{katz2017reluplex}, deciding injectivity or surjectivity \citep{froese2025parameterized,froese2024complexity} or deciding whether the Lipschitz constant of a ReLU network exceeds a certain threshold \citep{virmaux2018lipschitz,jordan2020exactly}.

Another line of research has studied the question of how to construct ReLU networks for functions with a certain number of regions \citep{he2018relu,chen2022improved,hertrich2023towards,brandenburg2025decomposition,zanotti2025linear}. The number of regions of maxout networks was studied by \citet{montufar2022sharp}. Note that all our hardness results hold for maxout networks, too, as maxout is a generalization of ReLU. \citet{goujon2024number} present bounds for general piecewise linear activation functions. The average number of linear regions was studied, among others, by \citet{hanin2019deep, tseran2021expected}. Our work is inspired by the aim to better understand complexity-theoretic aspects of neural networks; another well-studied question in that regime is the complexity of training \citep{goel2021tight,froese2022computational,froese2023training,bertschinger2023training}.

\section{Preliminaries}
\label{sec:prelim}
For $n\in \mathbb{N}$, we write $[n]:=\{1,\dots, n\}$.
For a set $P\subseteq \mathbb{R}^n$, we denote by $\overline{P}$, $P^\circ$, and $\partial P$ its closure, interior, and boundary, respectively.
The ReLU function is the real function $x\mapsto \max(0,x)$.
For any $n\in \mathbb{N}$, we denote by $\sigma:\mathbb{R}^n\to \mathbb{R}^n$ the function that computes the ReLU function in each component.

\paragraph{Polyhedra, CPWL functions, and hyperplane arrangements}
A \emph{polyhedron} $P$ is the intersection of finitely many closed halfspaces.
A \emph{polytope} is a bounded polyhedron.
A \emph{face} of $P$ is either the empty set or a set of the form $\arg\min \{c^\top x: x\in P\}$ for some $c\in \mathbb{R}^n$.
A \emph{polyhedral complex} $\mathcal{P}$ is a finite collection of polyhedra such that $\emptyset \in \mathcal{P}$, if $P\in \mathcal{P}$ then all faces of $P$ are in $\mathcal{P}$, and if $P,P'\in \mathcal{P}$, then $P\cap P'$ is a face of $P$ and $P'$.
A function $f:\mathbb{R}^n\to \mathbb{R}$ is \emph{continuous piecewise linear} (CPWL), if there exists a polyhedral complex $\mathcal{P}$ such that the restriction of $f$ to each full-dimensional polyhedron $P\in \mathcal{P}$ is an affine function.
If this condition is satisfied, then $f$ and $\mathcal{P}$ are \emph{compatible}.
A \emph{hyperplane arrangement} $\mathcal{H}$ is a collection of hyperplanes in $\mathbb{R}^n$.
A \emph{cell} of a hyperplane arrangement is an inclusion maximal connected subset of $\mathbb{R}^n\setminus(\bigcup_{H\in \mathcal{H}} H)$.
A hyperplane arrangement naturally induces an associated polyhedral complex with the cells being the maximal polyhedra of the complex.

\paragraph{ReLU networks}
A \emph{ReLU neural network} $N$ with $d\geq 0$ hidden layers is defined by $d+1$ affine transformations $T^{(i)}:\mathbb{R}^{n_{i-1}}\to \mathbb{R}^{n_i}, x\mapsto A^{(i)}x+b^{(i)}$ for $i\in [d+1]$.
We assume that $n_0=n$ and $n_{d+1}=1$.
The ReLU network $N$ \emph{computes} the CWPL function $f_N:\mathbb{R}^n\to \mathbb{R}$ with
\[
f_N=T^{(d+1)}\circ \sigma \circ \cdots \circ \sigma \circ T^{(1)}.
\]
The matrices $A^{(i)}\in \mathbb{R}^{n_i\times n_{i-1}}$ are called the \emph{weights} and the vectors $b^{(i)}\in \mathbb{R}^{n_i}$ are the \emph{biases} of the $i$-th layer.
We say the network has \emph{depth} $d+1$ and \emph{size} $s(N) := \sum_{i=1}^d n_i$.
Equivalently, ReLU networks can also be represented as layered, directed, acyclic graphs where each dimension of each layer is represented by one vertex, called a \emph{neuron}.
Each neuron computes an affine transformation of the outputs of its predecessors, applies the ReLU function, and outputs the result. We denote the CPWL function mapping the network input to the output of a neuron $v$ by $f_{N,v}:\mathbb{R}^n\to \mathbb{R}$. If the reference to the ReLU network $N$ is clear, we abbreviate $f_{N,v}$ by $f_v$.

\paragraph{Activation patterns}
Given a ReLU network $N$, a vector $a \in \{0,1\}^{s(N)}$ is called an \emph{activation pattern} of $N$ if there exists an input $x \in \mathbb{R}^n$ such that when $N$ receives $x$ as input, the $i$-th neuron in $N$ has positive output (is active) if $a_i=1$ and $0$ if $a_i=0$.
Given an activation pattern $a\in\{0,1\}^{s(N)}$, the network collapses to an affine function $f_N^a:\mathbb{R}^n\to \mathbb{R}$, and each neuron $i$ outputs an affine function $f^a_{N,i}:\mathbb{R}^n\to \mathbb{R}$ ($f^a_{N,i}$ is the zero function if $a_i=0$).
Again, if the reference to the ReLU network $N$ is clear, we abbreviate $f^a_{N,i}$ by $f^a_i$.

\paragraph{Encoding size}
We use $\langle \cdot \rangle$ to denote the encoding size of numbers, matrices, or entire neural networks, where we assume that numbers are integers or rationals encoded in binary such that they take logarithmic space. More details can be found in \Cref{sec:encoding}.

\paragraph{Computational Complexity}
We give an informal overview over some notions of computational complexity and refer to \citep{arora2009computational} for further reading.
A function $f:\{0,1\}^*\to \{0,1\}$ is in \P\ if $f$ is computable in polynomial time by a deterministic Turing machine, in \NP\ if it is computable in polynomial time by a non-deterministic Turing machine, and in \RP\ if it is computable in polynomial time by a randomized Turing machine that never outputs false positives and accepts a correct input with probability at least $1/2$.
Intuitively, \P\ contains problems that can be efficiently solved while \NP\ contains those whose solutions can be efficiently verified.
It widely believed that \P\,$\neq$\,\NP\ and \RP\,$\neq$\,\NP\ hold.
A function $f:\{0,1\}^*\to \mathbb{N}$ is in \#\P\ if there is a polynomial time non-deterministic Turing machine, which has exactly $f(x)$ accepting paths for any input $x\in \{0,1\}^*$ and in \FP\SPACE{} if $f$ is computable by a deterministic Turing machine that uses polynomial space.
A problem is called \emph{hard} for \NP\ (analogously, for \#\P) if all other problems in this class can be reduced to it in polynomial time, and \emph{complete} if it is both hard and contained in the class itself.

\section{Definitions of linear regions}\label{sec:definitions}

In this section we extract the six most commonly used definitions of linear regions from the literature and discuss their relations alongside with important properties and subtleties. \Cref{tab:literature} provides an overview of which previous papers use which definitions.

\begin{table}[b]
    \centering
    \hfill
    \begin{tabular}{ll}
        \toprule
         Paper & Definitions \\
         \midrule
         \citet{pascanu2013number} & \ref{def:region4} \\
         \citet{montufar2014number} & \ref{def:region5} \\
         \citet{raghu2017expressive} & \ref{def:region1}, \ref{def:region5} (*) \\
         \citet{arora2018understanding} & \ref{def:region5} \\
         \citet{serra2018bounding} & \ref{def:region1},\ref{def:region2} (*)\\
         \citet{hanin2019deep} & \ref{def:region2}, \ref{def:region4} \\
         \citet{he2018relu} & \ref{def:region3}, \ref{def:region6} \\
         \citet{rolnick2020reverse} & \ref{def:region2}, \ref{def:region5} (*) \\
         \citet{tseran2021expected} & \ref{def:region1},\ref{def:region5} (*)\\
         \citet{chen2022improved} & \ref{def:region3}, \ref{def:region5}, \ref{def:region6} \\
        \bottomrule
    \end{tabular}\hfill\hfill
    \begin{tabular}{ll}
        \toprule
         Paper & Definitions \\
         \midrule
         
         \citet{montufar2022sharp}& \ref{def:region4} \\
         \citet{wang2022estimation}& \ref{def:region1}, \ref{def:region5}  \\
         \citet{cai2023getting}& \ref{def:region2}  \\
         \citet{hertrich2023towards}& \ref{def:region6}  \\
         \citet{huchette2023deep}& \ref{def:region2} (*)  \\
         \citet{goujon2024number} & \ref{def:region3}, \ref{def:region6} \\
         \citet{brandenburg2025decomposition} & \ref{def:region3}, \ref{def:region6} \\
         \citet{masden2022algorithmic}& \ref{def:region2} \\
         \citet{zanotti2025bounds}& \ref{def:region4}, \ref{def:region6} \\
         \citet{zanotti2025linear} & \ref{def:region4} (*)\\
        \bottomrule
    \end{tabular}
    \hfill{}\vspace{.3em}
    
    \caption{List of papers that use one or several definitions.
    Additional notes on the papers marked with an asterisk can be found in \Cref{sec:additional_table_notes}. \citet{lezeau2024tropical} use another definition that lies between \Cref{def:region4,def:region5}, see \Cref{sec:flaws}.}
    \label{tab:literature}    
\end{table}

The set of inputs that have the same activation pattern induce a subset of $\mathbb{R}^n$ on which $f_N$ is affine.
\begin{definition}[Activation Region]
\label{def:region1}
    Given a network $N$ and an activation pattern $a\in\{0,1\}^{s(N)}$ with support $I\subseteq [s(N)]$, the set
    $
    S_{N,a}=\{x\in \mathbb{R}^n: f_i^a(x)>0\; \text{ for all } i\in I,\; f_i^a(x) \leq 0\; \text{ for all } i\notin I\}
    $
    is an \emph{activation region} of $N$. If the reference to the ReLU network is clear, we abbreviate $S_{N,a}$ by $S_a$.
\end{definition}
Activation regions can be open, closed, neither open nor closed, and full- or low-dimensional, see \Cref{fig:activation_regions} for some examples.
The (disjoint) union of all activation regions is exactly $\mathbb{R}^n$.
In particular, the number of activation regions equals the number of activation patterns.
It is important to note that the term \emph{activation region} is used ambiguously.
For example, \citet{hanin2019deep} use the term to refer to only \emph{full-dimensional} activation regions. 
\begin{definition}[Proper Activation Region]
\label{def:region2}
Given a ReLU network $N$, a \emph{proper activation region} of $N$ is a full-dimensional activation region of $N$.
\end{definition}

While the previous two definitions depend on the neural network \emph{representation} itself, the following four definitions depend only on the CPWL \emph{function} represented by the ReLU network and are independent from the concrete representation.

\begin{definition}[Convex Region]
\label{def:region3}
    Given a ReLU network $N$ and a polyhedral complex $\mathcal{P}$ that is compatible with $f_N$, a \emph{convex region of $N$ given $\mathcal{P}$} is a full-dimensional polyhedron $P\in \mathcal{P}$.
    The \emph{number of convex regions} of $N$ is the minimum number of convex regions of any polyhedral complex $\mathcal{P}$ that is compatible with $f_N$.
\end{definition}
Note that many different polyhedral complexes can attain this minimal number. Hence, in general, it is not possible to refer to a polyhedron $P$ as a `convex region of $N$' without specifying an associated polyhedral complex.

Another option to define linear regions is to use inclusion-maximal connected subsets on which the function computed by the ReLU network is affine, leading to the following definitions.
\begin{definition}[Open Connected Region]
\label{def:region4}
    Given a ReLU network $N$, an \emph{open connected region} of $N$ is an open, inclusion-wise maximal connected subset of $\mathbb{R}^n$ on which $f_N$ is affine.
\end{definition}
\begin{definition}[Closed Connected Region]
\label{def:region5}
    Given a ReLU network $N$, a \emph{closed connected region} of $N$ is a (closed) inclusion-wise maximal connected subset of $\mathbb{R}^n$ on which $f_N$ is affine.
\end{definition}

\begin{figure}[t]
  \centering
  \raisebox{-0.5\height}{
	\includegraphics[page=1, width=0.7\textwidth]{arxiv_figures.pdf}
    }
    \raisebox{-0.5\height}{
	\includegraphics[page=2, width=0.25\textwidth]{arxiv_figures.pdf}
    }
  \caption{
        A ReLU network computing the function $f(x,y)=\max(-y, \min(0,-x))$.
        The closed connected regions (center left) and the activation regions (center right) are displayed.
        The slice $\{(x,0):x\geq 0\}$ is contained in the two closed connected regions with functions $0$ (red) and $-y$ (blue).
		We have $R_6=R_5=R_4=3$, $R_3=4$, $R_2=5$ and $R_1=7$.
		In anti-clockwise direction starting from the region with value 0, the activation patterns are $010110$, $000000$, $001001$, $101001$, $100000$, $110010$ and $110110$ (neurons are ordered from the upper left to the lower right).}
		\label{fig:activation_regions}
\end{figure}

The subtle difference in the definition of open and closed connected regions has an important consequence:
As \citet{zanotti2025bounds} showed, $\overline{P_1}\cap\overline{P_2}=\partial P_1\cap \partial P_2$ holds for any distinct open connected regions $P_1,P_2$.
Interestingly, the same is \emph{not true} for closed connected regions. This is due to the fact that a closed connected region can continue on a low-dimensional slice of another closed connected region, which leads to a part of the boundary of one closed connected region to be contained in the interior of another closed connected region.
\citet[Figure 1]{zanotti2025bounds} gives a neat example where a low dimensional slice even connects two seemingly disconnected full-dimensional sets; another example can be found in \Cref{fig:activation_regions}.
Every open connected region is the interior of the closure of a union of some proper activation regions, see \Cref{lem:open_connected_region_is_made_up_of_proper_activation_regions}.
However, a closed connected region is in general not the closure of a union of some activation regions (see~\Cref{subsection:zanotti_network}).

\citet{hanin2019deep} define the set of open connected regions as the connected components of the input space where the set of points on which the gradient of $f_N$ is discontinuous are removed.
Alternatively, the set of open connected regions is equal to the unique set $\mathcal{S}$ with the minimum number of open connected subsets such that $\bigcup_{S\in \mathcal{S}}\overline{S}=\mathbb{R}^n$ and $f_N$ restricted to any $S\in \mathcal{S}$ is affine, see \Cref{lem:open_connected_regions_are_well_defined}.
The same is not true for closed connected regions, since there can be multiple sets $\mathcal{S}$ with the minimal number of closed connected subsets such that $\bigcup_{S\in \mathcal{S}}S=\mathbb{R}^n$ and $f_N$ restricted to any $S\in \mathcal{S}$ is affine. For example, in \Cref{fig:activation_regions}, in such a minimal set $\mathcal{S}$ there is exactly one closed subset corresponding to the closed connected region with the constant zero function. There are multiple options to choose this subset, e.g. $(-\infty, 0] \times [0,\infty)$ or $((-\infty, 0] \times [0,\infty)) \cup \{(x,0):x>0\}$.

By dropping the requirement of being connected, we obtain the following definition.
\begin{definition}[Affine Region]
\label{def:region6}
    Given a ReLU network $N$, an \emph{affine region} of $N$ is an inclusion-wise maximal subset of $\mathbb{R}^n$ on which $f_N$ is affine.
\end{definition}

For each definition, a linear region $S\subseteq \mathbb{R}^n$ of a ReLU network $N$ can be associated with an affine function $g:\mathbb{R}^n\to \mathbb{R}$ such that $f_N(x)=g(x)$ for all $x\in S$.
The affine function $g$ is unique if $S$ is full-dimensional.
We say that the function $g$ is \emph{computed} or \emph{realized} on $S$.
If $g$ is the zero function, we call $S$ a \emph{zero region} and a \emph{nonzero region} otherwise. The following theorem is immediate from the definitions.
\begin{theorem}
\label{thm:region_hierarchy}
    Given a ReLU network $N$, let $R_1, R_2,\dots,R_6$ denote the number of activation regions, proper activation regions, convex regions, open connected regions, closed connected regions and affine regions, respectively.
    Then:\quad
    $
    R_6 \leq R_5 \leq R_4 \leq R_3 \leq R_2 \leq R_1.
    $
\end{theorem}

The examples in \Cref{fig:activation_regions,fig:linear_regions} show that each inequality in \Cref{thm:region_hierarchy} can be strict.
\citet{zanotti2025bounds} showed $R_4\in O((R_6)^{n+1})$.
\citet{he2018relu} showed $R_3\leq (R_6)!$.
Trivially, $R_1\leq 2^{s(N)}$.

\begin{figure}
    \centering
    \hfill
    \includegraphics[page=4, width=0.3\textwidth]{arxiv_figures.pdf}\hfill
    \includegraphics[page=3, width=0.25\textwidth]{arxiv_figures.pdf}\hfill{}
    \caption{The function $\max(0,x)+\max(0,-y) - \max(0,x-y) +\min(\max(0, x-2), \max(0, y-2)) - 2 \min(\max(0, x-1), \max(0, y-1))$.
    We have $R_6=7$, $R_5=8$, and $R_4=9$.}
    \label{fig:linear_regions}
\end{figure}

\paragraph{Problem definitions for counting linear regions}
In the remainder of the paper, we consider algorithmic problems arising from counting linear regions. Both decision problems (e.g., deciding if the number of regions is larger than a given threshold) or function problems (such as computing exactly or approximately the number of linear regions) are detailed below.

\textsc{$K$-region-decision}
\\
\textbf{Input:} A ReLU network $N$.
\\
\textbf{Question:} Does $N$ have strictly more than $K$ linear regions (according to a specified definition)?

\textsc{Linear region counting}
\\
\textbf{Input:} A ReLU network $N$.
\\
\textbf{Question:} What is the number of linear regions of $N$ (according to a specified definition)?

\section{Counting regions: one hidden layer}
\label{sec:shallow}
In this section, we derive our results for ReLU networks with one hidden layer.
Our first main result is that \textsc{1-region-decision} can be solved in polynomial time for ReLU networks with one hidden layer.
Detailed proofs of the statements in this section are given in \Cref{sec:omitted_proofs_shallow}.
\begin{theorem}
\label{thm:polyresultfor1layerdecision}
\textsc{1-region-decision} for networks with one hidden layer is in \P{} for \Cref{def:region1,def:region2,def:region3,def:region4,def:region5,def:region6}.
\end{theorem}
The idea of the proof is as follows.
In a ReLU network $N$ with one hidden layer, each neuron corresponds to a hyperplane that divides the input space into two halfspaces.
It is not guaranteed that each hyperplane also leads to a discontinuity of the gradient of $f_N$, since the functions of the neurons with the same hyperplane may add up to an affine function.
The proof of \Cref{thm:polyresultfor1layerdecision} shows that detecting whether a hyperplane of a neuron is canceled can be done in polynomial time. All hyperplanes of the network cancel if and only if the network computes an affine function and has thus only one linear region according to \Cref{def:region3,def:region4,def:region5,def:region6}. For \Cref{def:region1,def:region2}, \textsc{1-region-decision} is trivial.

\citet[Lemma 15]{froese2024complexity} give a result similar to \Cref{thm:polyresultfor1layerdecision}.
They show that for a network with one hidden layer without biases, one can determine in polynomial time whether the network computes the constant zero function, and otherwise find a point on which the network computes a nonzero value.
In contrast, \Cref{thm:polyresultfor1layerdecision} considers biases and nonzero affine functions.

Turning to the problem of exactly counting the number of regions, we show the following theorem.
\begin{theorem}
\label{thm:linear_region_counting_is_sharp_p_complete_for_one_hidden_layer}
    \textsc{Linear region counting} for ReLU networks with one hidden layer is \#\P-hard for \Cref{def:region5,def:region6} and \#\P-complete for \Cref{def:region1,def:region2,def:region3,def:region4}.
\end{theorem}
\begin{proof}[Proof sketch]
Containment in \#\P{} is easy for \Cref{def:region1,def:region2}, since an activation pattern $a\in \{0,1\}^{s(N)}$ of a ReLU network $N$ is a unique certificate for a (proper) activation region, which can be verified in polynomial time by computing the dimension of the set $S_a$, see \Cref{lem:dim_of_activation_region_is_poly_time_computable}.
More modifications are necessary to show \#\P{} containment also for \Cref{def:region3,def:region4}.

To prove \#\P-hardness, we reduce from the problem of counting the number of cells of a hyperplane arrangement which is \#\P-complete, see \citep{linial1986hard}.
Starting from a hyperplane arrangement $\mathcal{H}$ in $\mathbb{R}^n$, we carefully construct a neural network whose linear regions exactly correspond to the cells of the hyperplane arrangement. With proper technical adjustments, this works for all six definitions.
\end{proof} 
\citet{linial1986hard} proved the \#\P-completeness of counting the number of cells of a hyperplane arrangement by reducing from the \#\P-complete problem of counting the number of acyclic orientations of a graph.
His reduction implies that \textsc{Linear region counting} remains \#\P-hard even for networks with one hidden layer where $A^{(2)}=(1,\dots,1)$, $b^{(1)}=0$, $b^{(2)}=0$, and $A^{(1)}$ is the transpose of an incidence matrix of a directed graph.

It is an open problem whether \textsc{Linear region counting} is in \#\P{} for \Cref{def:region5,def:region6}.
Notice that a single activation pattern does not suffice as a certificate, since two proper activation regions with a non-empty intersection can have the same affine function.
For example, consider the function $\max(0,x)+\max(0,-y)-\max(0,x-y)$ with zero regions $(-\infty, 0]\times [0,\infty)$ and $[0,\infty) \times (-\infty, 0]$.

To the best of our knowledge, it is unknown whether there is a polynomial factor approximation algorithm for approximating the number of cells in a hyperplane arrangement.
Thus, it is also an open problem whether \textsc{Linear region counting} has a polynomial factor approximation algorithm that runs in polynomial time.

\section{Counting regions: going beyond one hidden layer}
\label{sec:deep}
Here, we prove hardness results for ReLU networks with more than one hidden layer. Detailed proofs of the statements in this section are given in \Cref{sec:omitted_proofs_deep}, together with more detailed discussions providing additional intuition for some of the proofs.

\subsection{Hardness of the decision version}
From a result of \citet{wang2022estimation}, the following theorem follows immediately.
\begin{theorem}[\citep{wang2022estimation}]
    For any fixed constant $K\in \mathbb{N}_{\geq 1}$,\textsc{$K$-region-decision} for ReLU networks of depth $\Theta(\log n)$ is \NP-hard according to \Cref{def:region3,def:region4,def:region5,def:region6}.
\end{theorem}
In their reduction from 3-\SAT, they construct a network computing a minimum of $n+1$ terms. As known constructions for computing the minimum require depth $\Theta(\log n)$, this leads to hardness for counting regions of networks with depth $\Theta(\log n)$.
With \Cref{thm:np_hardness_Kregion_decision}, we improve on the result by showing that the problem remains $\NP$-hard even for networks with two hidden layers.
\begin{theorem}\label{thm:np_hardness_Kregion_decision}
    For any fixed constants $K, L\in \mathbb{N}_{\geq 1}$, $L\geq 2$, \textsc{$K$-region-decision} for ReLU networks with $L$ hidden layers is \NP-hard for \Cref{def:region3,def:region4,def:region5,def:region6}.
\end{theorem}

As a consequence, we even obtain hardness of the question whether there exists more than a single region. Proving this special case is also the first step of proving \Cref{thm:np_hardness_Kregion_decision}, as captured by the following lemma for the special case $K=1$ and $L=2$.
\begin{lemma}\label{thm:np_completeness_2region_decision_problem}
\textsc{1-region-decision} for ReLU networks with two hidden layers is \NP-complete according to \Cref{def:region3,def:region4,def:region5,def:region6}.
\end{lemma}
\begin{proof}[Proof sketch]
We reduce from \SAT. Given a \SAT{} formula $\phi$, we carefully construct a neural network $N_\phi$ with two hidden layers that has nonzero regions contained in $\varepsilon$-hypercubes around (0-1) points that satisfy $\phi$ and is constantly zero anywhere else. In this way, if $\phi$ is unsatisfiable, then $N_\phi$ computes the constant zero function and has exactly one linear region. If $\phi$ is satisfiable, then $N_\phi$ has strictly more than one linear region (one zero region and at least one nonzero region).
\end{proof}

Given a 3-SAT formula $\phi$ with $m$ clauses, the network $N_\phi$ from the reduction in the proof of \Cref{thm:np_completeness_2region_decision_problem} has input dimension and width $\mathcal{O}(m)$, whereas the network that is created in the reduction of \Cref{thm:np_hardness_Kregion_decision} has input dimension $\mathcal{O}(m)$ and width $\mathcal{O}(m+K)$.
We note that there is an alternative way to prove the \NP-hardness of \Cref{thm:np_completeness_2region_decision_problem}.
\citet[Theorem 18]{froese2024complexity} show that the problem of deciding whether or not a network without biases with one hidden layer has a point which evaluates to a positive value is \NP-complete.
By taking the maximum of the output of the network used in their reduction with the zero function, we obtain the \NP-hardness of \Cref{thm:np_completeness_2region_decision_problem}.
However, our reduction offers a new perspective on the difficulty of the problem. In fact, the ideas used in our reduction are built upon in \Cref{sec:hardness of approximation} to obtain results on the hardness of approximation of \textsc{Linear region counting}.
Moreover, our reduction has different properties, for example, all nonzero linear regions are bounded. This is not possible without biases, since then, all nonzero regions correspond to a union of polyhedral cones.

\Cref{thm:np_hardness_Kregion_decision} can be proven using \Cref{thm:np_completeness_2region_decision_problem} in two steps.
First, we can extend the hardness result of \Cref{thm:np_completeness_2region_decision_problem} from \textsc{1-region-decision} to \textsc{$K$-region-decision} by adding a new function with $K$ linear regions, and second, we can increase the number of hidden layers of the resulting network from $2$ to $L$ by adding $L-2$ additional hidden layers that compute the identity function.

As a corollary of \Cref{thm:np_hardness_Kregion_decision}, we obtain insights on the following decision problem.

\textsc{$L$-network-equivalence}
\\
\textbf{Input:} Two ReLU networks $N, N'$ with $L$ hidden layers.
\\
\textbf{Question:} Do the networks $N$ and $N'$ compute the same function?

Two ReLU networks compute the same function if and only if the difference of the networks is the zero function.
Since this difference can be computed by a single ReLU network, we obtain the following.
\begin{corollary}
\label{cor:euqality_of_two_relu_networks}
    \textsc{1-network-equivalence} is in \P{}, and, for any fixed constant $L\geq 2$, \textsc{$L$-network-equivalence} is $\coNP$-complete.
\end{corollary}
We also obtain the following runtime lower bound based on the Exponential Time Hypothesis.\footnote{The Exponential Time Hypothesis~\citep{impagliazzo2001complexity} states that 3-\SAT\ on $n$ variables cannot be solved in $2^{o(n)}$ time.}
\begin{corollary}
\label{cor:eth_lower_bound}
    For any fixed constants $K,L\in \mathbb{N}, L\geq 2$, \textsc{$K$-region-decision} and \textsc{Linear region counting} for \Cref{def:region3,def:region4,def:region5,def:region6} for ReLU networks with input dimension $n$ and $L$ hidden layers cannot be solved in $2^{o(n)}$ or $2^{o(\sqrt{\langle N\rangle})}$ time unless the Exponential Time Hypothesis fails.
\end{corollary}

The $2^{o(n)}$ lower bound can be seen as another example of the curse of dimensionality in machine learning.
As the input dimension grows, the problem quickly becomes intractable.

\subsection{Hardness of exact and approximate counting}
\label{sec:hardness of approximation}
Here, we show that even \emph{approximating} the number of linear region is hard for certain definitions.
We prove two inapproximability results for different network architectures.
For the first result, we use the proof ideas of \Cref{thm:np_completeness_2region_decision_problem} to show the following lemma.
\begin{lemma}
\label{thm:sharp_sat_reduction}
For any fixed constant $L\in \mathbb{N}, L\geq 3$, there is a reduction from $\#$\SAT\ to \textsc{Linear region counting} for networks with $L$ hidden layers according to \Cref{def:region4,def:region5}.
\end{lemma}
\begin{proof}[Proof sketch]
Given a \SAT\ formula $\phi$, the network $N_\phi$ from the proof of \Cref{thm:np_completeness_2region_decision_problem} has some nonzero linear regions contained in the $\varepsilon$-hypercube around every satisfying (0-1) point of $\phi$.
In order to get control over the number of linear regions created per satisfying point, we carefully need to modify the network $N_\phi$ such that every satisfying assignment of $\phi$ creates the same number of nonzero linear regions. This yields a simple formula relating the number of linear regions of the ReLU network with the number of satisfying assignments of $\phi$.
We achieve this by taking the minimum of the modified network with an appropriate function, which increases the number of hidden layers by one but does not change the width compared to $N_\phi$.
The reduction can be extended to ReLU networks with $L\geq 3$ hidden layers as before.
\end{proof}
We note that \Cref{thm:sharp_sat_reduction} does not hold for \Cref{def:region6}, since the constructed network has multiple closed connected regions with the same affine function.
\Cref{thm:sharp_sat_reduction} shows that a network having few regions does not necessarily imply that the regions of the network are ``easy to count''.
On the contrary, it shows that instances with relatively few regions can lead to \#\P-hard counting problems.

The reduction from \#\SAT{} implies that \emph{approximating} \textsc{Linear region counting} is intractable as well, in the following sense: We define an \emph{approximation algorithm} achieving approximation ratio $\rho \leq 1$ as an algorithm that is guaranteed\footnote{The algorithm may be probabilistic and return the correct answer with probability bounded away from $1/2$.} to return, given a network $N$ as input, a number that is at least $\rho$ times the number of regions of $N$. In fact, even though \Cref{thm:sharp_sat_reduction} only holds for \Cref{def:region4,def:region5}, it is sufficient to prove the inapproximability result also for \Cref{def:region3,def:region6}.
\begin{theorem}
\label{exponential_hardness_of_approximation_result}
    For any fixed constant $L\in \mathbb{N}, L\geq 3$, it is \NP-hard to approximate \textsc{Linear region counting} for \Cref{def:region3,def:region4,def:region5,def:region6} within an approximation ratio larger than $(2^n+1)^{-1}$ for networks with $L$ hidden layers and input dimension $n$.
\end{theorem}
\begin{proof}
If a \SAT\ formula has no satisfying assignment, the network produced by the reduction of \Cref{thm:sharp_sat_reduction} will have exactly $1$ linear region according to \Cref{def:region3,def:region4,def:region5,def:region6}. Otherwise, it will have at least $1+2^n$ linear regions according to \Cref{def:region3,def:region4,def:region5,def:region6}. If it was possible to achieve an approximation ratio larger than $(2^n+1)^{-1}$ in polynomial time we could decide if a \SAT\ formula is satisfiable in polynomial time. This concludes the proof.   
\end{proof}
We note that very similar ideas also rule out a fully polynomial randomized approximation scheme (FPRAS) for approximating the number of linear regions. A FPRAS \citep{jerrum2003counting} is a randomized polynomial-time (in the size of the input and $1/\varepsilon$) algorithm that returns a number $T$ such that $\text{Prob}[(1-\epsilon)R  \leq T\leq (1+\epsilon)R]\geq 3/4$,
where $R$ is the number of regions of the network. \Cref{exponential_hardness_of_approximation_result} can be easily adapted to show no FPRAS can exist for estimating the number of regions of neural network unless \NP\,=\,\RP.

\Cref{exponential_hardness_of_approximation_result} does not imply hardness of approximation for ReLU networks with two hidden layers.
In the following, we show that approximation is indeed hard for networks with two hidden layers for \Cref{def:region4,def:region5,def:region6}, although with a weaker inapproximability factor than in \Cref{exponential_hardness_of_approximation_result}.
\begin{theorem}
\label{thm:approx_result_2_hidden_layers}
Given a ReLU network with two hidden layers and input dimension $n$, for every $\varepsilon \in (0,1)$ it is \NP-hard to approximate \textsc{Linear region counting} by a ratio larger than $2^{-O(n^{1-\varepsilon})}$ for \Cref{def:region3,def:region4,def:region5,def:region6}.
\end{theorem}
\begin{proof}[Proof Sketch]
Given a \SAT{} formula $\phi$, the network $N_\phi$ in the proof of \Cref{thm:np_completeness_2region_decision_problem} has exactly one linear region if $\phi$ is unsatisfiable and at least 2 linear regions if $\phi$ is satisfiable.
We proceed by showing that for a ReLU network $N$ with $R$ linear regions, the network computing the function $f^{(k)}:\mathbb{R}^{nk}\to \mathbb{R},\; f^{(k)}(x_{11},...,x_{1n},...,x_{k1},...,x_{kn})=\sum_{i=1}^kf_{N}(x_{i1},...,x_{in})$
has exactly $R^k$ linear regions.
In words, $f^{(k)}$ is simply the sum of $k$ copies of $f_N$ each having as input a disjoint set of $n$ variables.
Applying this construction to $N_\phi$ for appropriate values of $k$ gives the desired result.
\end{proof}

\section{Counting regions using polynomial space}\label{sec:pspace}
Due to the \#\P-hardness of \textsc{Linear region counting}, we do not expect that efficient (polynomial time) algorithms for counting the number of linear regions exist.
\citet{wang2022estimation} claimed that \textsc{Linear region counting} for \Cref{def:region5} would be in \EXP\TIME. As stated in \Cref{sec:flaws}, the algorithm actually works for \Cref{def:region4} instead of \Cref{def:region5}. Since closed connected regions are generally not a union of activation regions, it is still an open problem whether even an \EXP\TIME{} algorithm is possible for \Cref{def:region5}. Also for \Cref{def:region3}, to the best of our knowledge, it is not clear whether an \EXP\TIME{} algorithm exists, because there are infinitely many options to choose the underlying polyhedral complex.
On the contrary, we show in this section that for \Cref{def:region1,def:region2,def:region6}, the number of regions can be computed in polynomial space and therefore also in \EXP\TIME{}.

\begin{theorem}
\label{thm:poly_space_algorithms}
    \textsc{Linear region counting} is in \FP\SPACE{} for \Cref{def:region1,def:region2,,def:region6}.
\end{theorem}
It is not hard to see that computing the number of activation regions and proper activation regions is possible in space that is polynomial in $\langle N\rangle$. Consider the following (informal) algorithm:
Given a ReLU network $N$, iterate over all $2^{s(N)}$ vectors in $\{0,1\}^{s(N)}$, and for each vector $a\in \{0,1\}^{s(N)}$, compute the dimension of $S_a$ in time polynomial in $\langle N\rangle$ (see~\Cref{lem:dim_of_activation_region_is_poly_time_computable}) to determine if $S_a$ is an activation region or a proper activation region, and increase a counter by one if this is the case.

Counting the number of affine regions is slightly more complicated, because a naive approach enumerating all the activation patterns would need to keep track of all the affine coefficients already seen to avoid double-counting, which is infeasible in polynomial space. Instead, we iterate over all possible affine functions by enumerating all possible coefficient combinations that have an encoding size of less than a polynomial upper bound, and check if there is a proper activation region on which the affine function is realized. The running time of this algorithm is exponential in the encoding size of the network, but it suffices to prove \FP\SPACE{} containment.

We describe how to count regions in \Cref{alg:exhaustive_search}.
The comments in the algorithms refer to the lemmas that show that the computation in the respective line can be performed in polynomial space.

\begin{algorithm}
\caption{\textsc{SearchAffinePiece}}
\label{alg:search_affine_piece}
    \begin{algorithmic}[1]
        \Require{A ReLU network $N$ and a vector $(a_1,\dots,a_n,b)\in \mathbb{Q}^{n+1}$.}
        \Ensure{1 if $\sum_{i=1}^n a_i x_i + b$ is a function of an affine region of $N$, else 0.}
        \For{$a\in \{0,1\}^{s(N)}$}
        \If{$\dim S_a=n$}\algorithmiccomment{(\Cref{lem:dim_of_activation_region_is_poly_time_computable})}
        \If{$\sum_{i=1}^n a_i x_i + b=f^a_{N}(x)$ }
        \Return 1 \algorithmiccomment{(\Cref{lem:pattern_to_function_is_poly_time_computable})}
        \EndIf
        \EndIf
        \EndFor
        \Return 0
    \end{algorithmic}
\end{algorithm}
\begin{algorithm}
\caption{\textsc{ExhaustiveSearch}}
\label{alg:exhaustive_search}
    \begin{algorithmic}[1]
        \Require{A ReLU network $N$.}
        \Ensure{Number of affine regions of $N$.}
        \State $n_{\max}=\max\{n_0,n_1,\dots, n_{d+1}\}$
        \State $U=2^{36d^2n_{\max}^2\langle A_{\max}\rangle}$ \algorithmiccomment{(\Cref{lem:coeffs_are_polynomial})}
        \State $R=0$
        \For{$(a, b)\in \{-U,\dots, U\}^{n+1}\times \{1,\dots, U\}^{n+1}$}
        \If{$\gcd(a_i,b_i)=1$ for $i\in [n+1]$}
        \State $R \leftarrow R+\textsc{SearchAffinePiece}(N,(\frac{a_1}{b_1},\dots, \frac{a_{n+1}}{b_{n+1}}))$ \algorithmiccomment{(\Cref{lem:checkaffine})}
        \EndIf
        \EndFor
        \Return $R$
    \end{algorithmic}
\end{algorithm}

The algorithms for counting (proper) activation regions (\Cref{def:region1,def:region2}) show fixed-parameter tractability of \textsc{Linear region counting} with respect to the number of neurons.
Furthermore, recent results of \citet[Corollary 4.4]{froese2025parameterized} imply W[1]-hardness of \textsc{Linear region counting} for ReLU networks with two or more hidden layers when parameterized by the input dimension (for \Cref{def:region3,def:region4,def:region5,def:region6}).
However, the fixed-parameter tractability status of \textsc{Linear region counting} remains open for other parameterizations and definitions.

\section{Conclusion}
We collected and discussed six commonly used non-equivalent definitions for linear regions of ReLU networks.
We proved \#\P-hardness for counting the number of linear regions (for all six definitions) and \NP-hardness for several associated decision problems (for most definitions).
We further showed that for ReLU networks with two or more hidden layers, even approximating the number of linear regions is \NP-hard (again, for most definitions).
On the positive side, we showed that for some definitions, linear regions can at least be counted in polynomial space.

There remain many interesting open problems and directions for future work.
Is \textsc{Linear region counting} in \EXP\TIME\ for \Cref{def:region3,def:region5} (are there even finite algorithms)?
Can some of the results in \Cref{sec:shallow,sec:deep,sec:pspace} be extended to all six definitions?
For example, is \textsc{Linear region counting} for ReLU networks with one hidden layer contained in \#\P\ also for \Cref{def:region5,def:region6}?
Is the problem of approximating the number of linear regions also \NP-hard for ReLU networks with one hidden layer?
Finally, it would be interesting to study the fixed-parameter tractability of \textsc{Linear region counting} under different parameterizations and definitions.

\bibliographystyle{plainnat}
\bibliography{arxiv.bib}


\appendix

\section{Notes on Theory and Literature}

\subsection{Details on the encoding size}
\label{sec:encoding}
The \emph{encoding size} $\langle n\rangle$ of a nonnegative integer $n$ is $\langle n\rangle:=\lceil\log_2(n+1)\rceil$, the encoding size of a fraction $q=a/b$ with $a\in \mathbb{Z}$ and $b\in \mathbb{N}$ is $\langle q\rangle := 1 + \langle |a|\rangle + \langle b\rangle$, and the encoding size of a matrix with rational entries $A=(a_{ij})_{i\in [n],\ j\in [m]}$ is $\langle A\rangle := nm + \sum_{i\in [n],\ j\in [m]}\langle a_{ij}\rangle$.
Since we are interested in the computational complexity, we restrict ourselves to ReLU networks with rational entries that can be represented with a finite number of bits.
Then, a ReLU network $N$ of depth $d+1$ with $A^{(i)}\in \mathbb{Q}^{n_i\times n_{i-1}}$ and a vector $b^{(i)}\in \mathbb{Q}^{n_i}$ has encoding size
$
\langle N\rangle = \sum_{i=1}^{d+1} \langle A^{(i)}\rangle + \sum_{i=1}^{d+1}\langle b^{(i)}\rangle
$.
In particular, if $A_{\max}$ denotes the maximum encoding size of an entry in any $A^{(i)}$ and $b^{(i)}$, then $\langle N\rangle \leq (d+1)\cdot (\max\{n,n_1,\dots,n_d\}+1)^2 \cdot (1+A_{\max})=\text{poly}(n,s(N),A_{\max})$.

\subsection{Basic properties of linear regions}
The following statements hold.
\begin{lemma}
\label{lem:open_connected_regions_are_well_defined}
    Given a ReLU network $N$, the set of open connected regions of $N$ is equal to the unique set $\mathcal{S}$ with the minimal number of open connected subsets such that $\bigcup_{S\in \mathcal{S}}\overline{S}=\mathbb{R}^n$ and $f_N$ restricted to any $S\in \mathcal{S}$ is affine.
\end{lemma}
\begin{proof}
    First, we show that the minimal set satisfying the assumptions above is unique.
    Suppose that there are two distinct sets $\mathcal{S}$ and $\mathcal{S}'$ of open subsets that achieve the minimum.
    By~\citet[Lemma 3.2]{zanotti2025bounds}, each element in $S\in \mathcal{S}$ is maximal in the sense that there does not exist a nonempty set $U\subseteq \mathbb{R}^n\setminus S$ such that $S\cup U$ is open and connected, and $f_N$ restricted to $S\cup U$ remains affine.
    
    Since $\mathcal{S}$ and $\mathcal{S}'$ are distinct, there is a set $S\in \mathcal{S}$ with $S\notin \mathcal{S}'$, and since $\bigcup_{S'\in \mathcal{S}'}\overline{S'}=\mathbb{R}^n$, there is a set $S'\in \mathcal{S}'$ with $S\cap S'\neq \emptyset$.
    Since both $S$ and $S'$ are open, the intersection $S\cap S'$ is full dimensional and the affine functions on $S$ and $S'$ are identical.
    Therefore, $U=S'\setminus S\subset \mathbb{R}^n \setminus S$ is a nonempty set such that $S\cup U=S\cup S'$ is open and $f_N$ restricted to $S\cup U$ remains affine, contradicting the maximality of $S$.

    By~\citet[Lemma 3.2]{zanotti2025bounds}, the unique set $\mathcal{S}$ is then exactly the unique set of inclusion-maximal open subsets of $\mathbb{R}^n$ such that $f_N$ restricted to each inclusion-maximal subset is affine, which is by definition the set of open connected regions of $N$.
\end{proof}
\begin{lemma}
\label{lem:open_connected_region_is_made_up_of_proper_activation_regions}
    Given a ReLU network, the closure of every open connected region is the closure of the union of a set of proper activation regions.
\end{lemma}
\begin{proof}
    See~\citet[Lemma 3]{hanin2019deep}.
\end{proof}
\subsection{Additional notes for definitions used in literature}
\label{sec:additional_table_notes}
\begin{itemize}
    \item \citet{raghu2017expressive} do not use \Cref{def:region1} explicitly, but they define activation patterns and derive a bound on the total number of activation patterns, which is equivalent to bounding the number of activation regions.
    \item The bound of \citet{serra2018bounding} holds for \Cref{def:region2}, compare the discussion in \citep[Section 5]{cai2023getting}. The MIP counts the number of activation regions (\Cref{def:region1}).
    \item \citet{rolnick2020reverse} only treat cases where \Cref{def:region2,def:region5} are equivalent.
    \item \citet{tseran2021expected} consider activation regions of maxout networks, which is conceptually slightly different from the activation regions defined in this paper.
    \item \citet{huchette2023deep} actually define activation regions using \Cref{def:region1}, but then later state that they disregard low dimensional linear regions. Therefore, their definition is equivalent to \Cref{def:region2}.
    \item \citet{zanotti2025linear} defines linear regions as the closure of open connected regions.
\end{itemize}
\subsection{Inaccuracies in the literature}
\label{sec:flaws}
In this section we list a few cases where misunderstandings about the different definition of a linear region led to small errors or inconsistencies in previous work, alongside with a suggestion how they would be fixable. Usually, this can be achieved by switching to a different, maybe more appropriate definition of a linear region.

\begin{enumerate}
    \item In Lemma 11 (d), \citep{chen2022improved} claims that $\left(S_1\cap S_2\right) \cap\left(S_1^\circ \cup S_2^\circ\right) = \emptyset$ holds for any two closed connected regions $S_1, S_2$. While this claim is not true for closed connected regions, it is true that $\left(\overline{S_1}\cap \overline{S_2}\right)\cap\left(S_1^\circ \cup S_2^\circ\right) = \emptyset$ holds for any two \emph{open} connected regions $S_1, S_2$. This inaccuracy was first pointed out by \citet{zanotti2025bounds}.
    
    \item The algorithm of \citet{wang2022estimation} does not create the set of closed connected regions, since in the algorithm, only closures of activation regions are merged such that (1) the affine function of the regions is identical and (2) the closures of the two activation regions have a non-empty intersection.
    However, a closed connected region is in general \emph{not} a union of the closure of a set of activation regions, see \Cref{subsection:zanotti_network} for an example.
    
    However, the algorithm of \citet{wang2022estimation} can instead be adapted to count the number of \emph{open connected regions} with one minor adjustment. Instead of merging the closure of activation regions with a nonempty intersection, merge only two proper activation regions if their closure has a $(n-1)$-dimensional intersection. This computation can still be done in time polynomial in the input size through linear programming, as described in \Cref{lem:dim_of_activation_region_is_poly_time_computable}.

    \item \citet{lezeau2024tropical} define a linear region as follows:
    
    \textit{A set of a neural network $f$ is a linear region if it is a maximal connected region (closure of an open set) on which $f$ is linear.}

    We note that their definition is not equivalent to \Cref{def:region4}~or~\ref{def:region5}, since a closed connected region that counts as multiple open connected regions can count as a single linear region in their definition, while a closed connected region can also count as multiple linear regions in their definition; consider, for example, the orange closed connected region in \Cref{fig:zanotti_network}.

    In Section 6.2, they present an algorithm to estimate the number of linear regions: they compute the number of unique gradients obtained on a set of sample points. They further state that if the gradients of two points are equal but the midpoint of the two points has a different gradient, then the two original points correspond to two different linear regions.
    This is not always true, since linear regions can be nonconvex according to their definition. Therefore, their algorithm can also overestimate the number of linear regions, but instead always underestimates the number of proper activation regions.
\end{enumerate}

\subsection{Technical results}
\begin{lemma}
\label{lem:coeffs_are_polynomial}
    Given a ReLU network $N$ and an activation pattern $a\in \{0,1\}^{s(N)}$ corresponding to a proper activation region. Let $A_{\max}\in \mathbb{Z}$ be the maximum absolute value of any numerator or denominator in an entry of the matrices $A^{(1)},\dots, A^{(d+1)}$ and biases $b^{(1)},\dots, b^{(d+1)}$ and let $n_{\max}=\max\{n_0,\dots,n_{d+1}\}$.
    Then, the encoding size of every coefficient of the affine function $f_N^a:\mathbb{R}^n\to \mathbb{R}$, $f_N^a(x)=\sum_{i=1}^n a_ix + b$ is bounded by
    \[
    36d^2n_{\max}^2\langle A_{\max}\rangle.
    \]
\end{lemma}
\begin{proof}
    Given an activation pattern $a\in \{0,1\}^{s(N)}$ of $N$, we modify the matrices $A^{(1)},\dots, A^{(d+1)}$ and biases $b^{(1)},\dots, b^{(d+1)}$ by replacing the columns of matrices and bias entries corresponding to inactive neurons with 0 entries.
    Let $A^{(1),a},\dots, A^{(d+1),a}$ and biases $b^{(1),a},\dots, b^{(d+1), a}$ denote the modified matrices and biases.
    Then, a simple calculation shows that for all $x\in \mathbb{R}^n$, we have
    \[
    f_N^a(x) = A^{(d+1),a}\cdot \dots \cdot A^{(1),a}x + b^{(d+1), a}+\sum_{i=0}^{d-1}A^{(d+1),a}\cdot \dots \cdot A^{(d+1-i),a}b^{(d-i), a}.
    \]
    To bound the encoding size of all occurring coefficients, we separately give a bound on the absolute value of any occurring denominator and numerator.
    For this, we turn the $d+1$ rational matrices and biases into integral matrices and biases by bringing all fractional entries to a common denominator.
    The value of the common denominator is bounded by $A_{\max}^{(d+1)(n_{\max}+1)^2}$, since there are fewer than $(d+1)(n_{\max}+1)^2$ entries in the rational matrices and biases.

    To bound the maximum absolute numerator value, we now consider the $(d+1)$ integral matrices and biases that arise by multiplying the fractional matrices by the common denominator.
    The maximum absolute value of an entry in one of the integral matrices or biases is bounded by $A_{\max}^{(d+1)(n_{\max}+1)^2}$.
    A simple calculation shows that the maximum absolute entry that can be obtained in the product of the $(d+1)$ integral matrices is
    \[
    \left(A_{\max}^{(d+1)(n_{\max}+1)^2}\right)^{d+1}\prod_{i=1}^{d}n_i
    \leq n_{\max}^d \cdot A_{\max}^{(d+1)^2(n_{\max}+1)^2}.
    \]
    The constant of $f_N^a$ is equal to $b^{(d+1), a}+\sum_{i=0}^{d-1}A^{(d+1),a}\cdot \dots \cdot A^{(d+1-i),a}b^{(d-i), a}$ and can thus be bounded by $(d+1) n_{\max}^d  A_{\max}^{(d+1)^2(n_{\max}+1)^2}$.
    
    Since the maximum encoding size of an element of a set of integers is obtained by the integer having the maximum absolute value, it follows that the encoding size of a coefficient in $f_N^a$ is at most
    \begin{align*}
    &1+\langle (d+1) n_{\max}^d  A_{\max}^{(d+1)^2(n_{\max}+1)^2} \rangle + \langle A_{\max}^{(d+1)(n_{\max}+1)^2} \rangle\\
    \leq\;&1+\langle d + 1 \rangle+d\langle n_{\max}\rangle+(d+1)^2(n_{\max}+1)^2\langle A_{\max}\rangle+(d+1)(n_{\max}+1)^2\langle A_{\max}\rangle\\
    \leq\;&1+\langle d + 1 \rangle+d\langle n_{\max}\rangle+2(d+1)^2(n_{\max}+1)^2\langle A_{\max}\rangle\\
    \leq\;&1+\langle d + 1 \rangle+d\langle n_{\max}\rangle+32d^2n_{\max}^2\langle A_{\max}\rangle\\
    \leq \; & 36d^2n_{\max}^2\langle A_{\max}\rangle.
    \end{align*}
\end{proof}
\begin{lemma}
\label{lem:pattern_to_function_is_poly_time_computable}
    Given a ReLU network $N$, an activation pattern $a\in \{0,1\}^{s(N)}$ and an index $i\in [s(N)]$ of a neuron, one can compute in time polynomial in $\langle N\rangle$ the (coefficients of the) affine function $f^a_i:\mathbb{R}^n\to \mathbb{R}$ which is computed at the output of the $i$-th neuron.
\end{lemma}
\begin{proof}
    The proof is analogous to the proof of \Cref{lem:coeffs_are_polynomial}.
\end{proof}
\begin{lemma}
\label{lem:dim_of_activation_region_is_poly_time_computable}
    Given a ReLU network $N$ and a vector $a\in \{0,1\}^{s(N)}$, one can compute the dimension of $S_a$ in time polynomial in $\langle N\rangle$.
\end{lemma}
\begin{proof}
    Let $I\subseteq [s(N)]$ denote the support of the activation pattern $a\in \{0,1\}^{s(N)}$.
    By \Cref{lem:coeffs_are_polynomial}, the encoding size of the coefficients in every affine function $f^a_i$ are bounded polynomially in $\langle N\rangle$. Thus, we can solve a series of linear programs to compute the dimension of the polyhedron
    \[
    P_a = \{x\in \mathbb{R}^n: f_i^a(x)\geq 0\; \text{ for all } i\in I,\; f_i^a(x) \leq 0\; \text{ for all } i\notin I\}.
    \]
    For details, we refer to~\cite[Section 7.3]{20.500.11850/426218}.
    Since $P_a$ is the closure of $S_a$ if $S_a$ is nonempty, it follows that the dimension of $P_a$ is equal to the dimension of $S_a$ unless $S_a$ is empty (in the latter case, $P_a$ cannot be full-dimensional).
    The latter case is easy to recognize since for all $i\in I$, we can check if $f_i^a(x)=0$ holds for all $x\in P_a$ by solving the linear program $\max\{f_i^a(x) : x\in P_a\}$.
\end{proof}
\section{Omitted proofs}
\subsection{Omitted proofs for the one hidden layer case}
\label{sec:omitted_proofs_shallow}
Let $N$ be a ReLU network with one hidden layer and let $[n_1]$ denote the set of neurons.
For ease of notation, we denote $A^{(2)}=(a_1,\dots,a_{n_1})\in \mathbb{R}^{1\times n_1}$ and $A^{(1)}=(w_{ij})_{i\in [n_1],j\in [n]}$.
Then, the function computed by the network is
\[
f_N(x)=b^{(2)}+\sum_{i=1}^{n_1} a_i \max(0, b^{(1)}_i + \sum_{j=1}^n w_{ij}x_j).
\]
For every neuron $i\in [n_1]$, we define the hyperplane $H_i := \{x\in \mathbb{R}^n : b^{(1)}_i + \sum_{j=1}^n w_{ij}x_j = 0\}$ and halfspaces $H_i^+ := \{x\in \mathbb{R}^n : b^{(1)}_i + \sum_{j=1}^n w_{ij}x_j \geq 0\}$, $H_i^- := \{x\in \mathbb{R}^n : b^{(1)}_i + \sum_{j=1}^n w_{ij}x_j \leq 0\}$.
The function $f_i:\mathbb{R}^n\mapsto \mathbb{R}$ computed by a neuron $i\in [n_1]$ is $f_i(x)=a_i \max(0, b^{(1)}_i + \sum_{j=1}^n w_{ij}x_j)$.
\begin{proof}[Proof of \Cref{thm:polyresultfor1layerdecision}]
We only consider \Cref{def:region3,def:region4,def:region5,def:region6} here (for \Cref{def:region1,def:region2}, the answer is trivially yes, as $n_1>0$).
Let $N$ be a ReLU network as defined above.

The idea is to check for every neuron if the function $f_N$ computed by the ReLU network $N$ has breakpoints along the hyperplane corresponding to the neuron, that is, we check if the gradient of $f_N$ is discontinuous along the hyperplane. Since multiple neurons can correspond to the same hyperplane, it is possible that the functions of these neurons cancel such that no breakpoints along the hyperplane are introduced. First, we group the neurons according to the hyperplanes to which they correspond.

Let $g_1(x)=b^{(1)}_1 + \sum_{j=1}^n w_{1j}x_j$ and let $I_1$ be the subset of $[n_1]$ such that
\[
i\in I_1 \quad \iff \quad c_{1i}= w_{i1} / w_{11}= \dots= w_{in} / w_{1n}= b^{(1)}_i / b^{(1)}_1
\]
holds for some constant $c_{1i}\in \mathbb{R}\setminus \{0\}$, which is equivalent to the hyperplanes $H_1,H_i$ being identical.

Now, we derive a condition when the function $f_N$ that is computed by the ReLU network has breakpoints along the hyperplane $H_1$.

For all $i\in I_1$, we have
\begin{equation*}
    f_i(x) = a_i \cdot \max(0, c_{1i}\cdot g_1(x)) = a_i |c_{1i}|\cdot \max(0, \sign(c_{1i})\cdot g_1(x)).
\end{equation*}

We split $I_1$ into the two sets $I_1^+=\{i\in I_1: c_{1i}>0\}$ and $I_1^-=\{i\in I_1: c_{1i}<0\}$.
The sum of the functions of the neurons in $I_1$ is
\begin{align*}
&\max(0, g_1(x))\sum_{i\in I_1^+}a_i c_{1i} - \max(0, -g_1(x))\sum_{i\in I_1^-}a_i c_{1i}\\
=\;&g_1(x)\sum_{i\in I_1^-}a_i c_{1i} + \max(0, g_1(x))\left(\sum_{i\in I_1^+}a_i c_{1i} - \sum_{i\in I_1^-}a_i c_{1i}\right)
\end{align*}
Thus, if 
\begin{equation}\label{eq:hyperplane_removal_condition}
\sum_{i\in I_1^+}a_i c_{1i} = \sum_{i\in I_1^-}a_i c_{1i},
\end{equation}
the sum of all functions having $H_1$ as corresponding hyperplane is affine and $f_N$ has no breakpoints along the hyperplane $H_1$.
Otherwise, $f_N$ has breakpoints along the hyperplane $H_1$ and $N$ has more than one linear region.

Thus, $N$ has only a single linear region if and only if \eqref{eq:hyperplane_removal_condition} holds for all $j\in[n_1]$ (replace 1 by $j$ in \eqref{eq:hyperplane_removal_condition} and define the set $I_j$ and constants $c_{ji}$ as before with $j$ instead of 1), which can be verified in polynomial time.
\end{proof}

\begin{proof}[Proof of \Cref{thm:linear_region_counting_is_sharp_p_complete_for_one_hidden_layer}]
We show \#P-hardness by reducing from the problem of counting the number of cells in a hyperplane arrangement, which is \#\P-complete (see~\citep{linial1986hard}).

Let $\mathcal{H}=(H_i)_{i\in [m]}$ be a hyperplane arrangement in $\mathbb{R}^n$ with $H_i:=\{x: w_i^\top x=0\}$, $w_i\in \mathbb{R}^n\setminus \{0\}$ such that each $w_i$ appears only once.
Restricted to such hyperplane arrangements, the problem of counting the number of cells remains \#\P-complete, see~\citep{linial1986hard}.
Given a point $x^*\in \mathbb{R}^n$ in a cell, we first orient the hyperplanes such that $w_i^\top x^* > 0$ for all $i\in [m]$.
We will show that the ReLU network $N_\mathcal{H}$ computing the convex function
\[
f_{N_\mathcal{H}}(x)=\sum_{i=1}^m \max(0, w_i^\top x),
\]
which can be computed using one hidden layer and one neuron per hyperplane, has exactly as many linear regions as the hyperplane arrangement $\mathcal{H}$ has cells (according to \Cref{def:region1,def:region2,def:region3,def:region4,def:region5,def:region6}).

It is easy to see that the number of cells of $\mathcal{H}$ is exactly the number of proper activation regions of $N_\mathcal{H}$.
We now show that also the number of activation regions is equal to the number of cells of $N_\mathcal{H}$.
Suppose for the sake of contradiction that there is an activation pattern $a\in \{0,1\}^m$ with support $I\subseteq [m]$ such that the activation region $S_a$ is neither full dimensional nor empty.
If $S_a$ is low dimensional, then by definition the set
\[
\{x\in \mathbb{R}^n: w_i^\top x \leq 0 \text{ for all }i\in [m]\setminus I\}
\]
is low dimensional and there is an index $j\in [m]\setminus I$ with $S_a\subseteq \{w_i^\top x = 0\}$.
Therefore, there is a $\lambda \in \mathbb{R}_{\geq 0}^{[m]\setminus I}$ such that
\[
\sum_{i\in [m]\setminus I} \lambda_i w_i = -w_j
\]
holds, which leads to the contradiction
\[
0<\sum_{i\in [m]\setminus I} \lambda_i w_i^\top x^* = -w_j^\top x^*\leq 0.
\]
Thus, each activation region is a proper activation region.

We now show that the affine functions on two cells cannot be equal, which proves \#\P-hardness also for \Cref{def:region3,def:region4,def:region5,def:region6}.

Suppose $a,a'\subseteq \{0,1\}^m$ are two activation patterns corresponding to distinct proper activation regions $S_{a}, S_{a'}$ with the same affine function.
$\overline{S_{a}}$ cannot have a $(n-1)$-dimensional intersection with $\overline{S_{a'}}$, since otherwise there would be exactly one hyperplane separating $S_{a}$ from $S_{a'}$, and by construction, the affine functions $f_{N_\mathcal{H}}^a$ and $f_{N_\mathcal{H}}^{a'}$ must be distinct.

Therefore, $\conv(S_a \cup S_{a'})\setminus (S_a \cup S_{a'})$ is full-dimensional, and there exists a proper activation region $S_{a^*}$ with $\dim(\text{conv}(S_{a}\cup S_{a'}) \cap S_{a^*})=n$ and $\dim(\overline{S_{a}}\cap \overline{S_{a^*}})=n-1$.
Since $f_{N_\mathcal{H}}$ is convex, $f_{N_\mathcal{H}}(x)=f_{N_\mathcal{H}}^a(x)$ holds for all $x\in \text{conv}(S_{a}\cup S_{a'})$. Thus, the function computed on $S_{a^*}$ must be equal to $f_{N_\mathcal{H}}^a$, which gives a contradiction as before.

We now show that \textsc{Linear region counting} is in \#\P\ for \Cref{def:region1,def:region2,def:region3,def:region4}.

A certificate for \Cref{def:region1,def:region2} is simply an activation pattern $a\in \{0,1\}^{s(N)}$ of the ReLU network $N$, which can be checked in polynomial time by computing the dimension of $S_a$, see \Cref{lem:dim_of_activation_region_is_poly_time_computable}.
Thus, \textsc{Linear region counting} is in \#\P\ for \Cref{def:region1,def:region2}.

We now construct certificates for \Cref{def:region3,def:region4}.
Given a ReLU network $N$ with one hidden layer, the set $\mathcal{H}=(H_i)_{i\in[m]}$ of hyperplanes that correspond to breakpoints of the function $f_N$ can be computed in polynomial time using the procedure described in the proof of \Cref{thm:polyresultfor1layerdecision}.
By construction, the function $f_N$ is affine on each cell of the hyperplane arrangement $\mathcal{H}$ and the affine functions that are realized on two neighboring cells (cells with an $(n-1)$-dimensional intersection) cannot be equal.
Thus, each cell of the hyperplane arrangement $\mathcal{H}$ is an open connected region.

Since each open connected region is convex, the set of convex regions of the ReLU network $N$ is well defined and its cardinality is equal to the set of open connected regions.

A unique certificate for an open connected region (and a convex region) is now given by the hyperplane arrangement $\mathcal{H}$ as well as a vector $a\in\{-,+\}^m$ specifying a cell $C\subset\mathbb{R}^n$ of the hyperplane arrangement $\mathcal{H}$, where $C\subseteq H_i^{a_i}$ for every $i\in [m]$.
The certificate can be checked in polynomial time: we can verify in polynomial time if the hyperplane arrangement $\mathcal{H}$ is defined as above, and we can verify in polynomial time if the vector $a$ corresponds to a cell of $\mathcal{H}$.

It follows that \textsc{Linear region counting} is in \#\P{} for \Cref{def:region3,def:region4}.
\end{proof}

\subsection{Omitted proofs for more than one hidden layers}
\label{sec:omitted_proofs_deep}
\textbf{Intuition for the proof of \Cref{thm:np_completeness_2region_decision_problem}.}
The proof of \Cref{thm:np_completeness_2region_decision_problem} improves the reduction of~\citet{wang2022estimation}, who use a result of \citet[Appendix I]{katz2017reluplex}.
Similar to \citet[Appendix I]{katz2017reluplex}, we rely on the simple fact that given a \SAT{} formula $\phi(x)=\bigwedge_{i=1}^m((\bigvee_{j\in J^+_i}x_j) \vee (\bigvee_{j\in J^-_i}\neg x_j))$ on the variables $x_1,\dots, x_n$, the function
\[
g_\phi(x) = 1 - \sum_{i=1}^m \max(0, 1-\sum_{j\in J^+_i}x_j-\sum_{j\in J^-_i}(1-x_j))
\]
takes value 1 on all satisfying (0-1) assignments, and value less than 0 on all non-satisfying assignments, which follows from the fact that $\max(0, 1-\sum_{j\in J^+_i}x_j-\sum_{j\in J^-_i}(1-x_j))$ evaluates to 0 for all (0-1) assignments that satisfy the $i$-th clause of $\phi$, and to 0 for all (0-1) assignments that do not satisfy the $i$-th clause of $\phi$.
For every $i \in [m]$, $J_i^+$ and $J^-_i$ are disjoint subsets of $[n]$ specifying which (negated) variables occur in the $i$-th clause of $\varphi$.

Notice that if $\phi$ is unsatisfiable, there is no (0-1) assignment on which $g_\phi$ takes value 1. As a result, for any $\varepsilon\in (0,1)$, $\phi$ is satisfiable if and only if $\max(1, \varepsilon + g_\phi) - 1 = \max(0, \varepsilon-1 + g_\phi)$ evaluates to $\varepsilon$ on some (0-1) assignment.

This implies that if $\phi$ is satisfiable, the function $h_{\phi, \varepsilon}= \max(0, \varepsilon-1 + g_\phi)$ has at least two linear regions according to \Cref{def:region3,def:region4,def:region5,def:region6}, since $h_{\phi, \varepsilon}$ evaluates to $\varepsilon$ for a satisfying (0-1) point, and to 0 for all points in an $\varepsilon$-ball around a non-satisfying (0-1) point. Since each clause in a \SAT{} formula is not satisfied by least one (0-1) assignment, we can assume that $\phi$ has a non-satisfying assignment.

If for any \SAT{} formula $\phi$, the function $h_{\phi, \varepsilon}$ had strictly more than one linear region (according to \Cref{def:region3,def:region4,def:region5,def:region6}) only if $\phi$ is satisfiable, then we would have a complete reduction from \SAT{} to the problem of deciding whether a ReLU network with two hidden layers has strictly more than one linear region (according to \Cref{def:region3,def:region4,def:region5,def:region6}), since $h_{\phi, \varepsilon}$ can be computed using a ReLU network with two hidden layers.

Unfortunately, there exists a \SAT{} formula $\psi$ such that $h_{\psi, \varepsilon}$ has more than one linear region although $\psi$ is unsatisfiable, see \Cref{ex:satfornphardness2}.

The key idea to resolve this is to add a CWPL function that is negative everywhere but on the elements of the set $\{0,1\}^n$ (on which it evaluates to zero).

A function with this property is the function $T_n:\mathbb{R}^n\to \mathbb{R}$ with
\[
T_n(x)=\sum_{i=1}^n(-\max(0,-x_i)-\max(0,x_i)+\max(0,2x_i-1)-\max(0,2x_i-2)),
\]
shown in \Cref{fig:subtraction_function3}. The proof of \Cref{thm:np_completeness_2region_decision_problem} shows that adding this function recovers the equivalence: A \SAT{} formula is satisfiable if and only if the function $\max(0, T_n + \varepsilon-1 + g_\phi)$ has more than one linear region according to \Cref{def:region3,def:region4,def:region5,def:region6}.

For an example that visualizes the different steps of the reduction, see \Cref{ex:satfornphardness1}.

\begin{figure}
    \centering
    \includegraphics[page=5, width=0.3\textwidth]{arxiv_figures.pdf}
    \includegraphics[page=6, width=0.325\textwidth]{arxiv_figures.pdf}
    \includegraphics[page=7, width=0.275\textwidth]{arxiv_figures.pdf}
    \caption{The function $T_1$ (left), $T_2$ (center) and its linear regions (right).}
    \label{fig:subtraction_function3}
\end{figure}

The following lemma is easy to prove based on the plot of $T_1$ in \Cref{fig:subtraction_function3}.
\begin{lemma}
\label{lem:t_regions}
For any $n\in \mathbb{N}$, the following implications hold
\begin{align*}
     &T_n(x)\leq -\varepsilon
     \quad&&\Longleftarrow&&\quad
     \exists i\in [n]:\;x_i\in (-\infty,-\varepsilon]\cup[\varepsilon,1-\varepsilon]\cup [1+\varepsilon,\infty)\\
     &T_n(x)=0
     \quad&&\iff&&\quad
     x\in \{0,1\}^n.
    \end{align*}
\end{lemma}

\begin{proof}[Proof of \Cref{thm:np_completeness_2region_decision_problem}]
We first show that the problem is in \NP.
If a ReLU network $N$ is a yes-instance of \textsc{1-region-decision}, then there are two proper activation regions on which two distinct affine functions are computed.
Two activation patterns corresponding to such proper activation regions serve as a polynomial certificate of a yes-instance.
Given two vectors $a,a'\in \{0,1\}^{s(N)}$, we can verify the certificate in polynomial time.
First, we check if $a$ and $a'$ correspond to proper activation regions by computing the dimension of $S_a$ and $S_{a'}$ in polynomial time, see \Cref{lem:dim_of_activation_region_is_poly_time_computable}.
If $S_a$ and $S_{a'}$ are proper activation regions and $f_N^a\neq f_N^{a'}$ holds, which can be checked in polynomial time (see \Cref{lem:pattern_to_function_is_poly_time_computable}), then $a,a'$ is a valid certificate of a yes-instance.
Thus, the problem is in \NP.

To show \NP-hardness, we reduce the problem of deciding whether a \SAT\ instance is satisfiable to our problem.
Let $\phi(x)=\bigwedge_{i=1}^m((\bigvee_{j\in J^+_i}x_j) \vee (\bigvee_{j\in J^-_i}\neg x_j) $ be a \SAT\ formula on $n$ variables with $|J^+_i|+|J^-_i|\leq n$.
Set $\varepsilon = 1/(n+1)$.
Consider the network $N_\phi$ with two hidden layers that computes the function
\[
    f_{N_\phi}(x)=\max(0, T_n(x)+\varepsilon-\sum_{i=1}^m \max(0, 1-\sum_{j\in J^+_i}x_j-\sum_{j\in J^-_i}(1-x_j))),
\]
Note that $N_\phi$ can be constructed from $\phi$ in polynomial time, since adding $T_n$ increases the encoding size only by an additional $\mathcal{O}(n^2)$ term.
The idea is now to show that if $\phi$ has a satisfying assignment, then $N_\phi$ has at least two linear regions, and if $\phi$ has no satisfying assignment, then $N_\phi$ has only one linear region with the constant zero function, which proves the lemma.

By \Cref{lem:t_regions}, we have $T_n(x)\leq -\varepsilon$ and therefore $f_{N_\phi}(x)=0$ for all $x\in 
\mathbb{R}^n$ with some $x_i\in (-\infty,-\varepsilon]\cup[\varepsilon,1-\varepsilon]\cup [1+\varepsilon,\infty)$.

As a result, we have
\[
\{x: f_{N_\phi}(x)>0\}\subseteq ([-\varepsilon,\varepsilon]\cup [1-\varepsilon, 1+\varepsilon])^n= \bigcup_{x\in \{0,1\}^n}B^\infty_\varepsilon(x),
\]
where $B^\infty_\varepsilon(x):=\{x': \|x-x'\|_\infty \leq \varepsilon\}$.

Suppose now that $x^*\in \{0,1\}^n$ satisfies $\phi$.
Then,
\[
\sum_{j\in J^+_i}x^*_j+\sum_{j\in J^-_i}(1-x^*_j)\geq 1 \quad \text{ for all }i\in [m],
\]
which implies $f_{N_\phi}(x^*)=\max(0, T_n(x^*)+\varepsilon)=\varepsilon>0$.
Thus, $N_\phi$ has at least two linear regions.

Suppose now that $x^*\in \{0,1\}^n$ does not satisfy $\phi$.
There is at least one clause $i^*$ with
\[
\sum_{j\in J^+_{i^*}}x^*_j+\sum_{j\in J^-_{i^*}}(1-x^*_j)= 0.
\]
In particular, for all $x\in B^\infty_\varepsilon(x^*)$, we have
\[
1-\sum_{j\in J^+_{i^*}}x_j-\sum_{j\in J^-_{i^*}}(1-x_j)
\geq
1 - |J^+_{i^*}|\varepsilon - |J^-_{i^*}|\varepsilon
\geq 1 - n \varepsilon
= 1 - n / (n+1)
= 1 / (n+1)
=\varepsilon.
\]
Therefore, we have $f_{N_\phi}(x)=0$ for all $x\in B^\infty_\varepsilon(x^*)$.
If $\phi$ has no satisfying assignment, then $f_{N_\phi}$ is the constant zero function.
\end{proof}
\begin{proof}[Proof of \Cref{thm:np_hardness_Kregion_decision}]
Given fixed constants $K, L\in \mathbb{N}_{\geq 1}$, $L\geq 2$ and a \SAT{} formula $\phi$, we will create a network with $L$ hidden layers which has strictly more than $K$ linear regions if and only if the network $N_\phi$ from the proof of \Cref{thm:np_completeness_2region_decision_problem} has strictly more than one hidden layer.

Let $N_\phi$ be the network as in the proof of \Cref{thm:np_completeness_2region_decision_problem}.
If $K\geq 2$, the network $N^{(K)}_\phi$ computing the function
\[
f_{N_\phi}(x)
-
\max(0,2(n+m)(x_1-2))-\dots - \max(0,2(n+m)(x_1-K))
\]
has $K$ linear regions if $N_\phi$ has only one linear region and strictly more than $K$ linear regions if $N_\phi$ has more than one linear region.
For \Cref{def:region3,def:region4,def:region5}, this follows from the fact that the newly introduced linear regions are outside of the hypercube $[-\varepsilon,1+\varepsilon]^n$ that contains all nonzero linear regions of $f_{N_\phi}(x)$.
For \Cref{def:region6}, we additionally have to verify that no newly introduced affine function was already present in $f_{N_\phi}$.
To see that no newly affine function was already present in $f_{N_\phi}$, observe that the coefficient of $x_1$ of every newly introduced affine function is at most $-2(n+m)$, while the coefficient of every affine function of $f_{N_\phi}$ cannot be smaller than $-n-m$, which can be easily seen from the proof of \Cref{thm:np_completeness_2region_decision_problem}.

The additional maximum terms can be created using two neurons in the first hidden layer that correspond to the positive and negative part of $x_1$, respectively, and adding $K-1$ neurons in the second hidden layer (using the equation $x_1=\max(0, x_1)-\max(0,-x_1)$ to build the maximum terms in the second hidden layer).

Thus, the network $N^{(K)}_\phi$ has two hidden layers.
To obtain a network $N^{(K,L)}_\phi$ with $L$ hidden layers, we add $L-2$ new hidden layers between the output layer and the last hidden layer of $N^{(K)}_\phi$.
Each new hidden layer has two neurons, the first neuron outputs $\max(0, N^{(K)}_\phi)$ and the second neuron outputs $\max(0, -N^{(K)}_\phi)$.
We achieve this by replacing the arcs from the second hidden layer of $N^{(K)}_\phi$ to the output node by connections to the two neurons of the first newly added hidden layer.
The theorem now follows by noting that the modified network $N^{(K,L)}_\phi$ has encoding size $\mathcal{O}(K\cdot\langle N_\phi\rangle+L)$ and can be constructed from $N_\phi$ in polynomial time.
\end{proof}
\begin{proof}[Proof of \Cref{cor:euqality_of_two_relu_networks}]
    Let $L\in \mathbb{N}$ be a fixed constant.
    Given two ReLU networks $N, N'$ with $L$ hidden layers, let $N^-$ represent the network with $L$ hidden layers that `subtracts' $N'$ from $N$ by computing the networks $N$ and $N'$ in parallel. $N$ and $N'$ compute the same function if and only if $N^-$ computes the zero function.

    To see that \textsc{$L$-network-equivalence} is in \NP, note that a vector in $\{0,1\}^{s(N^-)}$ that corresponds to a proper activation pattern with a nonzero affine function can be used as a certificate, as in the proof of \Cref{thm:np_completeness_2region_decision_problem}.

    If $L=1$, by \Cref{thm:polyresultfor1layerdecision} we can decide in polynomial time if $N^-$ computes an affine function.
    If $N^-$ computes an affine function then it computes the zero function if and only if $N^-$ evaluates to zero on $n+1$ affinely independent points, which yields a polynomial time algorithm for \textsc{1-network-equivalence}.
    Suppose $L\geq 2$ and let $\phi$ be a \SAT{} formula, let $N^{(1)}_\phi$ be the ReLU network with $L$ hidden layers from the proof of \Cref{thm:np_hardness_Kregion_decision}, and let $N_0$ be a ReLU network with $L$ hidden layers that computes the zero function.
    By \Cref{thm:np_hardness_Kregion_decision}, a \SAT{} formula $\phi$ is satisfiable if and only if $N^{(1)}_\phi$ and $N_0$ are a no-instance of \textsc{$L$-network-equivalence}, proving that \textsc{$L$-network-equivalence} is \coNP-hard.
\end{proof}

\begin{proof}[Proof of \Cref{cor:eth_lower_bound}]
Let $K, L\in \mathbb{N}, L\geq 2$ be fixed constants.
Given a 3-\SAT\ formula $\phi$ on $n$ variables and $m$ clauses, let $N^{(K,L)}_\phi$ be the ReLU network with $L$ hidden layers and input dimension $n$ from the proof of
\Cref{thm:np_hardness_Kregion_decision}. Recall that $N^{(K,L)}_\phi$ has strictly more than $K\in \mathbb{N}$ hidden layers if and only if $\phi$ is satisfiable.

We now show that $N^{(K,L)}_\phi$ has encoding size $\mathcal{O}(m^2)$. Recall that $N^{(K,L)}_\phi$ has an encoding size of $\mathcal{O}(K\cdot \langle N_\phi\rangle+L)$, and $N_\phi$ has an encoding size of $\mathcal{O}(n^2+nm)$. Since $n\leq 3m$ holds for every 3-\SAT\ formula and $K$ and $L$ are constants, the encoding size of $N^{(K,L)}_\phi$ is $\mathcal{O}(m^2)$.

It is well known that, assuming the Exponential Time Hypothesis is true, this implies that there is no $2^{o(n)}$ or $2^{o(\sqrt{\langle N\rangle})}$ time algorithm for \textsc{$K$-region-decision}, see~\citep{cygan2015lower}.
A $2^{o(n)}$ or $2^{o(\sqrt{\langle N\rangle})}$ time algorithm for \textsc{Linear region counting} problem would directly give a $2^{o(n)}$ or $2^{o(\sqrt{\langle N\rangle})}$ time algorithm for \textsc{$K$-region-decision}.
\end{proof}

\textbf{Intuition for the proof of \Cref{thm:sharp_sat_reduction}}

Given a \SAT\ formula $\phi$, the network $N_\phi$ from the proof of \Cref{thm:np_completeness_2region_decision_problem} has some nonzero linear regions near every satisfying assignment of $\phi$.
Unfortunately, the number of linear regions created per satisfying point depends on the formula $\phi$ and is not easily computable.
Therefore, we modify the network $N_\phi$ such that the same number of nonzero linear regions is created by every satisfying assignment of $\phi$. For this, we take the minimum of $f_{N_\phi}$ with the function
$T_{n,\varepsilon}:\mathbb{R}^n\to \mathbb{R}$,
\[
T_{n,\varepsilon}(x)= \max(0, \varepsilon+T_n(x)),
\]
shown in \Cref{fig:subtraction_function4}. We further show that $T_{n,\varepsilon}$ is strictly smaller than $f_{N_\phi}$, but greater than zero near a satisfying point. In this way, the minimum of $f_{N_\phi}$ and $T_{n,\varepsilon}$ is attained by $f_{N_\phi}$ near each non-satisfying point (where $f_{N_\phi}$ equals the zero function) and by $T_{n,\varepsilon}$ near each satisfying point. We proceed by showing that exactly $2^n$ nonzero regions are created for every satisfying assignment of $\phi$, and if $\phi$ has exactly $k$ satisfying assignments, the modified network has exactly $1+k\cdot 2^n$ linear regions according to \Cref{def:region4,def:region5}.
The modified network has three hidden layers. The reduction can then be extended to ReLU networks with $L\geq 3$ hidden layers as before.

\begin{figure}
    \centering
    \includegraphics[page=8, width=0.3\textwidth]{arxiv_figures.pdf}
    \includegraphics[page=9, width=0.325\textwidth]{arxiv_figures.pdf}
    \includegraphics[page=10, width=0.275\textwidth]{arxiv_figures.pdf}
    \caption{The functions $T_{1,\varepsilon}$ (left), $T_{2,\varepsilon}$ (center) and its linear regions (right).}
    \label{fig:subtraction_function4}
\end{figure}

The following lemma is required for the proof of \Cref{thm:sharp_sat_reduction}.
\begin{lemma}
\label{lem:teps_regions}
Let $n\geq 2$ and $B^\infty_\varepsilon(x):=\{x'\in \mathbb{R}^n: \|x-x'\|_\infty \leq \varepsilon\}$.
The function $T_{n,\varepsilon}$ with $0<\varepsilon< 1/2$
has exactly $1+2^{2n}$ linear regions according to \Cref{def:region4,def:region5} and we have
\begin{align*}
    &T_{n,\varepsilon}(x)>0&&\quad\Longrightarrow\quad&& x\in\bigcup_{x^*\in \{0,1\}^n}B^\infty_\varepsilon(x^*),\\
    &T_{n,\varepsilon}(x)=\varepsilon&&\quad\iff\quad&&x\in\{0,1\}^n,
\end{align*}
and for every $x^*\in \{0,1\}^n$, the set $B^\infty_\varepsilon(x^*)$ contains exactly $2^n$ nonzero regions according to \Cref{def:region3,def:region4,def:region5,def:region6}.
\end{lemma}
\begin{proof}[Proof of \Cref{lem:teps_regions}]
By \Cref{lem:t_regions}, if $T_n(x)\geq -\varepsilon$, then $x\in B^\infty_\varepsilon(x^*)$ for some $x^*\in \{0,1\}^n$.
Therefore, we have $T_{n,\varepsilon}(x) = 0$ for all $x\in \mathbb{R}^n\setminus (\bigcup_{x^*\in \{0,1\}^n}B^\infty_\varepsilon(x^*))$.
What is left is to analyze the linear regions of $T_{n,\varepsilon}$ in the set $B^\infty_\varepsilon(x^*)$ for every $x^*\in \{0,1\}^n$.

Due to the symmetry of $T_{n,\varepsilon}$, we only consider the set $B^\infty_\varepsilon((1,\dots, 1)^\top)=[1-\varepsilon,1+\varepsilon]^n$.
We will show that $[1-\varepsilon,1+\varepsilon]^n$ has exactly $2^n$ nonzero linear regions.

First, observe that for a point $x\in [1-\varepsilon,1+\varepsilon]^n$, we have
\[
T_n(x)= \sum_{i: x_i< 1}(x_i-1)+\sum_{i: x_i>1}(1-x_i).
\]
Given a subset $I$ of $[n]$, we define the set
\[
C_I:=\{x: 1-\varepsilon\leq x_i< 1 \;\forall i\in I,\;
1\leq x_i\leq 1+\varepsilon\;\forall i\notin I\}.
\]
It is easy to see that the disjoint union $\bigcup_{I\subseteq [n]}C_I$ gives exactly the set $[1-\varepsilon,1+\varepsilon]^n$.

Each set $C_I$ divides into two sets:
\begin{align*}
    &C^1_I:=\{x\in C_I: \sum_{i\in I}(x_i-1)+\sum_{i\notin I}(1-x_i)\geq-\varepsilon\}\\
    &C^0_I:=\{x\in C_I: \sum_{i\in I}(x_i-1)+\sum_{i\notin I}(1-x_i)\leq-\varepsilon\}
\end{align*}
We have $T_{n,\varepsilon}(x)=\varepsilon+\sum_{i\in I}(x_i-1)+\sum_{i\notin I}(1-x_i)$ for all $x\in C^1_I$ and $T_{n,\varepsilon}(x)=0$ for all $x\in C^0_I$.

The set $C^1_I$ is full dimensional, as $x^*\in \mathbb{R}^n$ with
$x^*_i=\begin{cases}
1-\frac{\varepsilon}{2n},&i\in I\\
1+\frac{\varepsilon}{2n},&i\notin I
\end{cases}$ is an interior point of $C^1_I$.
This proves that every $C_I$ contains exactly one nonzero region. Since the function for every $C_I$ is unique, $[1-\varepsilon,1+\varepsilon]^n$ contains exactly $2^n$ nonzero regions according to \Cref{def:region3,def:region6}.
Since $T_{n,\varepsilon}$ has only a single zero region, it follows that $T_{n,\varepsilon}$ has exactly $1+2^{2n}$ linear regions.
\end{proof}
\begin{proof}[Proof of \Cref{thm:sharp_sat_reduction}]
Let $\phi(x)=\bigwedge_{i=1}^m((\bigvee_{j\in J^+_i}x_j) \vee (\bigvee_{j\in J^-_i}\neg x_j) $ be a \SAT\ formula on $n$ variables, where $|J^+_i|+|J^-_i|\leq n$.
Set $\varepsilon = 1/(2+n+nm)$.
Consider the network $N_\phi^*$ that computes the function
\[
    f_{N_\phi^*}(x)=\min(T_{n,\varepsilon}(x), \max(0, 1-(n+1)\varepsilon-\sum_{i=1}^m \max(0, 1-\sum_{j\in J^+_i}x_j-\sum_{j\in J^-_i}(1-x_j)))),
\]
which can be computed with 3 hidden layers. This is due to the fact that the minimum of two terms can be expressed using three neurons $\min(a,b)=-\max(0,b-a)+\max(0,b)-\max(0,-b)$.
The reduction is polynomial since the addition of $T_{n,\varepsilon}$ increases the encoding size only by an additional $\mathcal{O}(n^2)$ term.
Now, our goal is to show that if $\phi$ has exactly $k$ satisfying assignments, then $N_\phi^*$ has exactly $1+2^n \cdot k$ linear regions.

By \Cref{lem:teps_regions}, if $f_{N_\phi^*}(x)> 0$ then $x\in \bigcup_{x^*\in \{0,1\}^n}B^\infty_\varepsilon(x^*)$, where $B^\infty_\varepsilon(x^*)=\{x'\in \mathbb{R}^n: \|x^*-x'\|_\infty \leq \varepsilon\}$.

We will prove our theorem by showing that the following holds for all $x^*\in \{0,1\}^n$.
\begin{enumerate}
    \item If $\phi(x^*)=0$, then $f_{N_\phi^*}$ has no nonzero linear region in $B^\infty_\varepsilon(x^*)$.
    \item If $\phi(x^*)=1$, then $f_{N_\phi^*}$ has exactly $2^n$ nonzero linear regions in $B^\infty_\varepsilon(x^*)$.
\end{enumerate}
We start with the first implication.
Suppose $\phi(x^*)=0$ holds.
Then, there is a clause $i^*$ such that
$x^*_j= \begin{cases}
    0,\; j\in J^+_{i^*}\\
    1,\;j\in J^-_{i^*}
\end{cases}$
holds.
Thus, we have for all $x\in B^\infty_\varepsilon(x^*)$:
\begin{align*}
&-\sum_{i=1}^m \max(0, 1-\sum_{j\in J^+_i}x_j-\sum_{j\in J^-_i}(1-x_j))\\
\leq &-\max(0,1-\sum_{j\in J^+_{i^*}}x_j-\sum_{j\in J^-_{i^*}}(1-x_j))\\
\leq &-\max(0,1-\sum_{j\in J^+_{i^*}}\varepsilon-\sum_{j\in J^-_{i^*}}(1-(1-\varepsilon)))\\
=&-1+(|J^+_{i^*}|+|J^-_{i^*}|) \varepsilon,
\end{align*}
implying
\begin{align*}
1-(n+1)\varepsilon-\sum_{i=1}^m \max(0, 1-\sum_{j\in J^+_i}x_j-\sum_{j\in J^-_i}(1-x_j))
\leq (|J^+_{i^*}|+|J^-_{i^*}|-n-1)\varepsilon
\leq 0
\end{align*}
and thus $f_{N_\phi^*}(x)=0$ for all $x\in B^\infty_\varepsilon(x^*)$.

To prove the second implication, suppose that $\phi(x^*)=1$ holds.
We will show that the second component $f_{N_\phi}$ in the minimum of $f_{N_\phi^*}$,
\[
f_{N_\phi}(x)=\max(0, 1-(n+1)\varepsilon-\sum_{i=1}^m \max(0, 1-\sum_{j\in J^+_i}x_j-\sum_{j\in J^-_i}(1-x_j)))
\]
is greater or equal to $T_{n,\varepsilon}$ for all $x\in B^\infty_\varepsilon(x^*)$. Then $f_{N_\phi^*}(x)=T_{n,\varepsilon}(x)$ holds for all $x\in B^\infty_\varepsilon(x^*)$. By \Cref{lem:teps_regions}, this implies that $f_{N_\phi^*}$ has exactly $2^n$ nonzero linear regions in $B^\infty_\varepsilon(x^*)$, which will prove the second implication.

We now show that $f_{N_\phi^*}(x)=T_{n,\varepsilon}(x)$ holds for all $x\in B^\infty_\varepsilon(x^*)$.
W.l.o.g. let $x^*=(1,\dots, 1)^\top$.
By assumption, $|J^+_i|\geq 1$ and $|J^-_i|\leq n-1$ holds for all clauses $i\in[m]$.
Thus, we have for all $x\in B^\infty_\varepsilon(x^*)=[1-\varepsilon,1+\varepsilon]^n$ and all $i\in [m]$
\begin{align*}
    1-\sum_{j\in J^+_i}x_j-\sum_{j\in J^-_i}(1-x_j)
    \leq 1-|J^+_i|(1-\varepsilon)+|J^-_i|\varepsilon
    \leq  \varepsilon+(n-1)\varepsilon
    = n\cdot \varepsilon.
\end{align*}
As a consequence, for all $x\in [1-\varepsilon,1+\varepsilon]^n$, we have
\begin{align*}
    f_{N_\phi}(x)
    &\geq \max(0, 1-(n+1)\varepsilon-\sum_{i=1}^m \max(0, 1-\sum_{j\in J^+_i}x_j-\sum_{j\in J^-_i}(1-x_j)))\\
    &\geq \max(0, 1-(n+1)\varepsilon-\sum_{i=1}^m n\cdot \varepsilon)\\
    &= 1-(n+1)\varepsilon-m \cdot n\cdot \varepsilon\\
    &= 1-(1+n+nm)\varepsilon\\
    &= \varepsilon\\
    &\geq T_{n,\varepsilon}(x),
\end{align*}
and thus, $f_{N_\phi^*}(x)=T_{n,\varepsilon}(x)$ for all $x\in [1-\varepsilon,1+\varepsilon]^n$.
By \Cref{lem:teps_regions}, $T_{n,\varepsilon}$ has $2^n$ nonzero linear regions in $[1-\varepsilon,1+\varepsilon]^n$.

We extend the hardness result to networks with $L\geq 3$ hidden layers as in the proof of \Cref{thm:np_hardness_Kregion_decision}.
\end{proof}
The following lemma will be used in the proof of \Cref{thm:approx_result_2_hidden_layers}.
\begin{lemma}
\label{lem:amplification_lemma}
Let $g:\mathbb{R}^n\to \mathbb{R}$ be a CWPL function with exactly $m$ affine regions. Then, for every $k\in \mathbb{N}$, the function $g^{(k)}:\mathbb{R}^{nk}\to \mathbb{R}^n$,
\[
g^{(k)}(x_{1,1},\dots, x_{1,n},\dots,x_{k,1},\dots, x_{k,n})=\sum_{i=1}^k g(x_{i,1},\dots, x_{i,n})
\]
has exactly $m^k$ affine regions.
\end{lemma}
\begin{proof}
Let $U_1,\dots,U_m$ be the affine regions of $g$, let $R_1,\dots, R_p$ be the affine regions of $g^{(k)}$ and let $h_i:\mathbb{R}^{n}\to \mathbb{R}$ be the affine function of the affine region $U_i$ of $g$ for every $i\in [m]$.

For every $i\in [m]^k$ and for all $x\in U_{i_1}\times \dots \times U_{i_k}$, the function $g^{(k)}$ computes the affine function
\[
g^{(k)}(x_{1,1},\dots, x_{1,n},\dots,x_{k,1},\dots, x_{k,n})=\sum_{j=1}^k h_{i_j}(x_{j,1},\dots, x_{j,n}).
\]
Since all affine functions $h_1,\dots,h_m$ are distinct, $U_{i_1}\times \dots \times U_{i_k}$ is contained in a different affine region of $g^{(k)}$ for every $i\in [m]^k$.
As $U_{i_1}\times \dots \times U_{i_k}$ is inclusion-maximal with respect to affinity of $g^{(k)}$, it follows that $\{R_1,\dots,R_p\} = \{U_{i_1}\times \dots \times U_{i_k}: i\in [m]^k\}$, which concludes the proof.
\end{proof}
\begin{proof}[Proof of \Cref{thm:approx_result_2_hidden_layers}]
Let $\phi$ be a \SAT{} formula on $l$ variables and let $N_\phi$ be the ReLU network with two hidden layers constructed in the proof of \Cref{thm:np_completeness_2region_decision_problem}. Recall that for \Cref{def:region3,def:region4,def:region5,def:region6}, the network $N_\phi$ has at least two linear regions if $\phi$ is satisfiable and exactly one linear region if $\phi$ is unsatisfiable.

Let $N_\phi^{(k)}$ be the ReLU network composed of taking $k$ copies of $N$ each with a disjoint set of $l$ variables.
The function computed by the ReLU network $N^{(k)}$ is $f_{N_\phi^{(k)}}:\mathbb{R}^{lk}\to \mathbb{R}$ with
\[
f_{N_\phi^{(k)}}(x_{11},\dots, x_{1l},\dots,x_{k1},\dots, x_{kl})=\sum_{i=1}^k f_{N_\phi}(x_{i1},\dots, x_{il}).
\]
If $\phi$ is unsatisfiable then $N_\phi$ has exactly one linear region which implies that $N_\phi^{(k)}$ also has exactly one linear region (according to \Cref{def:region3,def:region4,def:region5,def:region6}). If $\phi$ is satisfiable then $N_\phi$ has at least two affine regions and by \Cref{lem:amplification_lemma}, $N_\phi^{(k)}$ has at least $2^k$ affine regions. By \Cref{thm:region_hierarchy}, $N_\phi^{(k)}$ then also has at least $2^k$ linear regions according to \Cref{def:region3,def:region4,def:region5}.
It follows that approximating the number of regions within a factor larger than $2^{-k}$ is \NP-hard (according to \Cref{def:region3,def:region4,def:region5,def:region6}). Setting $n=lk$, the theorem now follows by picking $k=l^C$ for a sufficiently large constant $C$ (e.g., $C$ such that $\frac{C}{C+1}>1-\varepsilon$) and noting the construction of $N_\phi^{(k)}$ from $N_\phi$ can be done in polynomial time.
\end{proof}

\subsection{Omitted proofs for polynomial space algorithms}
\begin{lemma}\label{lem:checkaffine}
Given a ReLU network $N$ and an affine map $\varphi(x_1,...,x_n)=\sum_{i=1}^n a_ix_i+b$,
we can check in space polynomial in $\langle N\rangle$ and in the encoding size of the coefficients of $\varphi$ whether $\varphi$ is the function of an affine region of $N$.
\end{lemma}
\begin{proof}
First, note that $\varphi$ is the function of an affine region of $N$ if and only if there is a proper activation region on which $\varphi$ is realized.

Now, go over all $2^{s(N)}$ possible activation patterns for neurons of $N$ using space polynomial in $s(N)$. For each vector $a\in\{0,1\}^{s(n)}$, it is possible to verify in time polynomial in the encoding size $\langle N \rangle$ of the ReLU network $N$ whether $S_a$ is a proper activation region, see \Cref{lem:dim_of_activation_region_is_poly_time_computable}.
Further, we can check in polynomial time whether $\varphi$ is equal to the function $f^a_N$ computed on the proper activation region $S_a$, see \Cref{lem:pattern_to_function_is_poly_time_computable}.
\end{proof}
\begin{proof}[Proof of \Cref{thm:poly_space_algorithms}]
Let $N$ be a ReLU network.
As discussed in \Cref{sec:pspace}, the number of activation regions and proper activation regions can be counted in space which is polynomial in the encoding size $\langle N \rangle$ of the ReLU network $N$.
Now, we describe a polynomial space algorithm for counting the number of affine regions.

By \Cref{lem:coeffs_are_polynomial}, the encoding size of any coefficient of an affine function that occurs in one of the affine regions of $N$ is bounded by $M:=36d^2n_{\max}^2\langle A_{\max}\rangle$, which is polynomial in $\langle N \rangle$.

As each coefficient of an affine function of $N$ is a fraction, $M$ is also an upper bound on the encoding size of a numerator and on the encoding size of a denominator.
Since each affine function of $N$ is defined by $n+1$ fractions, we can exhaustively search through all sequences of $n+1$ fractions, where the numerator and denominator of each fraction can have encoding size of at most $M$. For each sequence of $n+1$ fractions, by \Cref{lem:checkaffine} we can compute in space that is polynomial in $\langle N \rangle$ if the corresponding affine function is the function of an affine region of $N$.
If an affine function of an affine region is found, we increase a counter by 1.
To avoid counting the same affine function more than once, we only check fraction sequences win which the numerator and denominator of every fraction are relatively prime.
\end{proof}
\section{Examples}
\subsection{A closed connected region which is not a closure of a union of a set of activation regions}
\label{subsection:zanotti_network}
\citet[Figure 1]{zanotti2025bounds} uses the following function as an example:
\[
\min(y, \max(-1, -x), \max(3 - 2 x, -x)).
\]
We turn this function into a ReLU network $N$ with three hidden layers, as illustrated in \Cref{fig:zanotti_network}.
For the orange closed connected region $P$ in \Cref{fig:zanotti_network}, there is no set of activation regions such that $P$ is the closure of a union of activation regions of the ReLU network $N$.

\begin{figure}
    \centering
    \includegraphics[page=17, width=0.45\textwidth]{arxiv_figures.pdf}
    \includegraphics[page=18, width=0.35\textwidth]{arxiv_figures.pdf}
    \caption{A ReLU network $N$ computing $\min(y, \max(-1, -x), \max(3 - 2 x, -x))$.
    An activation region of $N$ is either a blue line, blue point, or a full dimensional cell as defined by the blue lines.
    There are four closed connected region as indicated by the colors. The line between the points $(1,-1)$ and $(3,-3)$ belongs to the green as well as the orange region.}
    \label{fig:zanotti_network}
\end{figure}

\subsection{Further examples}
\begin{example}
\label{ex:satfornphardness1}
Consider the \SAT\ formula $\varphi = (\neg x_1)\land (x_1 \lor x_2)$ with the satisfying assignment $(0,1)$ and the function $g_\varphi(x)=1-\max(0,1-(1-x_1))-\max(0,1-x_1-x_2)$ displayed in \Cref{fig:sat_example_figure}.

\begin{figure}
    \centering
    \includegraphics[page=11, width=0.34\textwidth]{arxiv_figures.pdf}
    \includegraphics[page=12, width=0.34\textwidth]{arxiv_figures.pdf}
    \includegraphics[page=13, width=0.34\textwidth]{arxiv_figures.pdf}
    \includegraphics[page=14, width=0.34\textwidth]{arxiv_figures.pdf}
    \includegraphics[page=15, width=0.34\textwidth]{arxiv_figures.pdf}
    \includegraphics[page=16, width=0.34\textwidth]{arxiv_figures.pdf}
    \caption{The functions $g_\varphi(x)=1-\max(0,x_1)-\max(0,1-x_1-x_2)$ (top), $h_{\varphi, \varepsilon}(x)=\max(0, \varepsilon-1+g_\varepsilon(x))$ (center), and $f_{N_\varphi}(x)=\max(0,T_2(x)+\varepsilon-1+g_\varphi(x))$ (bottom) for $\varepsilon=0.4$.
    The function $f_{N_\varphi}$ is only nonzero in the blue region, which is contained in the $\varepsilon$-square (orange) around the only satisfying point of $\varphi$ (black).
    }
    \label{fig:sat_example_figure}
\end{figure}

As mentioned above, for all $x\in \{0,1\}^2$, $g_\varphi(x)=1$ holds if $x$ is a satisfying assignment of $\varphi$ and $g_\varphi(x)\leq 0$ otherwise.
Since $\varphi$ has an satisfying assignment, the function $h_{\varphi, \varepsilon}$ with $h_{\varphi, \varepsilon}(x)=\max(0, \varepsilon-1+g_\varphi(x))$ has strictly more than one linear region, see \Cref{fig:sat_example_figure}.
The final function in the reduction of is $f_{N_\varphi}(x)=\max(0, T_2(x)+\varepsilon-1+g_\varphi(x))$, see \Cref{fig:sat_example_figure}.
\end{example}
\begin{example}
\label{ex:satfornphardness2}
Consider the \SAT\ formula and function
\begin{align*}
\psi=\;&(x_1\lor x_2)\land
(\neg x_1\lor x_2)\land
(x_1\lor \neg x_2)\land
(\neg x_1\lor \neg x_2),\\
h_{\psi, \varepsilon}(x_1,x_2)=\;& \max(0, \varepsilon-\max(0,1-x_1-x_2)-\max(0,1-(1-x_1)-x_2)\\
&\phantom{\ \max(0, \varepsilon}-\max(0,1-x_1-(1-x_2))-\max(0,1-(1-x_1)-(1-x_2))).
\end{align*}
It is clear that $\psi$ is unsatisfiable.
However, for every $\varepsilon> 0$ we have $h_{\psi, \varepsilon}(1/2,1/2)=\varepsilon>0$.
\end{example}


\end{document}

\typeout{get arXiv to do 4 passes: Label(s) may have changed. Rerun}
